\documentclass[12pt, a4paper, twoside]{book}

%%%%%%%%%%%Pacchetti Per Compilazione%%%%%%%%%%%
\usepackage[utf8]{inputenc}
\usepackage{graphicx}
\usepackage[a-1b]{pdfx}
\usepackage[justification=centering]{caption}
\usepackage{subcaption}
\usepackage{fancyhdr} %abstract
\usepackage[italian]{babel}
\usepackage{listings}
\usepackage{amsthm}
\usepackage{amsmath}
\usepackage{todonotes}
\usepackage{color}
\usepackage{algorithm}
\usepackage{frontespizio}
\usepackage[noend]{algpseudocode}
\usepackage{tabularx}
\usepackage{hhline}
\usepackage[section]{placeins}

\usepackage{blindtext} % Per Dummy Text - remove
%%%%%%%%%%%%%%%%%%%%%%%%%%%%%%%%%

%%%%%%%%%%%Comandi per formattazione%%%%%%%%%%%
%Per mantenere i floats (immagini, tabelle etc.) all'interno delle sottosezioni
\makeatletter
\AtBeginDocument{%
	\expandafter\renewcommand\expandafter\subsection\expandafter{%
		\expandafter\@fb@secFB\subsection
	}%
}
\makeatother

%Label in italiano
\makeatletter
\renewcommand{\ALG@name}{Algoritmo}
\makeatother

\renewcommand{\lstlistingname}{Codice}

\renewcommand{\listalgorithmname}{Elenco degli algoritmi}

%Pagina bianca post capitolo effettivamente vuota e senza header
\let\origdoublepage\cleardoublepage
\newcommand{\clearemptydoublepage}{%
	\clearpage
	{\pagestyle{empty}\origdoublepage}%
}
\let\cleardoublepage\clearemptydoublepage

%%%%%%%%%%%%%%%%%%%%%%%%%%%%%%%%%

%%%%%%%%%%%Gestione Codice in Listings (Per C++)
\definecolor{mygreen}{rgb}{0,0.4,0}
\definecolor{mygray}{rgb}{0.35,0.35,0.35}
\definecolor{mymauve}{rgb}{0.58,0,0.82}

\lstset{ %
	backgroundcolor=\color{white},   % choose the background color; you must add \usepackage{color} or \usepackage{xcolor}
	basicstyle=\footnotesize,        % the size of the fonts that are used for the code
	breakatwhitespace=false,         % sets if automatic breaks should only happen at whitespace
	breaklines=true,                 % sets automatic line breaking
	captionpos=b,                    % sets the caption-position to bottom
	commentstyle=\color{mygray},    % comment style
	deletekeywords={...},            % if you want to delete keywords from the given language
	escapeinside={(*@}{@*)},          % if you want to add LaTeX within your code
	extendedchars=true,              % lets you use non-ASCII characters; for 8-bits encodings only, does not work with UTF-8
	frame=no,	                 % adds a frame around the code
	keepspaces=true,                 % keeps spaces in text, useful for keeping indentation of code (possibly needs columns=flexible)
	keywordstyle=\color{mygreen},       % keyword style
	language=C++,                    % the language of the code
	otherkeywords={*,...},           % if you want to add more keywords to the set
	numbers=left,                    % where to put the line-numbers; possible values are (none, left, right)
	numbersep=5pt,                   % how far the line-numbers are from the code
	numberstyle=\tiny\color{black}, % the style that is used for the line-numbers
	rulecolor=\color{black},         % if not set, the frame-color may be changed on line-breaks within not-black text (e.g. comments (green here))
	showspaces=false,                % show spaces everywhere adding particular underscores; it overrides 'showstringspaces'
	showstringspaces=false,          % underline spaces within strings only
	showtabs=false,                  % show tabs within strings adding particular underscores
	stepnumber=1,                    % the step between two line-numbers. If it's 1, each line will be numbered
	stringstyle=\color{black},     % string literal style
	tabsize=2,	                   % sets default tabsize to 2 spaces
	title=\lstname % show the filename of files included with \lstinputlisting; also try caption instead of title
}


%%%Macro per pseudocodice
\newcommand*\Let[2]{\State #1 $\gets$ #2}
%%%%%%%%%%%%%%%%%%%%%%


%%%%%%%%%%%Gestione Headers
\newcommand{\fncyblank}{\fancyhf{}}
\pagestyle{fancy}
\fancyhead{}
\fancyhead[LE, RO]{\leftmark}

%%%%%%%%%%%Alcuni Environment Utili
\theoremstyle{plain}
\newtheorem{thm}{Teorema}[section]
\newtheorem{cor}[thm]{Corollario}
\newtheorem{lem}[thm]{Lemma}
\newtheorem{prop}[thm]{Proposizione}

\theoremstyle{definition}
\newtheorem{defn}{Definizione}[section]
\newtheorem{example}{Esempio}[section]

\theoremstyle{remark}
\newtheorem{oss}{Osservazione}

\begin{document}

\selectlanguage{italian}
%%%% Il frontespizio
%Per modificare:
% 1) Modifica Frontespizio/Frontespizio.tex
% 2) Compila Tesi (questo file)
% 3) Compila Tesi-frn.tex
% 4) Ricompila Tesi (questo file)
\begin{frontespizio}
	\Margini{3cm}{3cm}{3cm}{3cm}
	\Istituzione{}
	\Logo[6cm]{img/logo_unipr}
	\Dipartimento{Scienze Matematiche, Fisiche e Informatiche}
	\Corso[Laurea]{Informatica}
	\Annoaccademico{2022/2023}
	\Titoletto{Tesi di Laurea}
	\Titolo {\uppercase {Realizzazione di un plugin per l'API Gateway Kong per l'integrazione con un microservizio per il controllo degli accessi tramite token}}
	\Candidato{Valerio Desiati}
	\Relatore{Prof. Roberto Alfieri}
\end{frontespizio}


%%%% La dedica
\newpage
\null\vspace{\stretch{1}}
\begin{flushright}
	\textit{Dedica}
\end{flushright}
\vspace{\stretch{3}}\null
\newpage

%%%% Gli indici
\pagestyle{plain}
\pagenumbering{roman}
\tableofcontents
%
\listoffigures    %Commentare se non vi sono Immagini
\listofalgorithms %Commentare se non vi sono Algoritmi
\listoftables     %Commentare se non vi sono Tabelle
%
%
%
%%%% La prefazione
\chapter*{Introduzione} %Se si cambia il Titolo cambiare anche la riga successiva così che appia corretto nell'indice
\addcontentsline{toc}{chapter}{Introduzione} %Per far apparire Introduzione nell'indice (Il nome deve rispecchiare quello del chapter)
\pagenumbering{arabic} % Settaggio numerazione normale
% L'introduzione deve contenere un riassunto del lavoro di Tesi.
% In particolare bisogna esprimere chiaramente e molto sinteticamente: contesto dello studio, motivazioni, contributo e conclusioni.
% Bisogna quindi fare un sommario dello studio ad alto livello, fornendo le intuizioni senza ricadere in dettagli tecnici.
% Anche lo stile dovrebbe essere più discorsivo rispetto alle parti tecniche della tesi.

Prima di addentrarsi nelle specifiche di questo progetto è bene partire dal concetto di applicazioni cloud – native.\\
Un'applicazione cloud – native è un'applicazione concepita e realizzata per risiedere in cloud.\\
L'approccio da utilizzare per lo sviluppo di questa tipologia di applicazioni è diametralmente opposto a quello per lo sviluppo di un'applicazione monolitica, 
come sarà spiegato in seguito nel paragrafo \ref{sec:microserviziintro}.\\ \\
Le applicazioni cloud – native si basano su tre concetti fondamentali:
\begin{itemize}
\item Orchestratori di container;
\item Microservizi;
\item Scalabilità.
\end{itemize}

Lo scopo di questo progetto è stato di realizzare un plugin per l'API Gateway Kong.\\
Il plugin prevede l'integrazione di Kong con un microservizio per verificare l'avvenuto acquisto di un modulo applicativo da parte di un utente, identificato tramite token.
Il progetto includerà la realizzazione del microservizio e l'ottimizzazione del plugin tramite l'utilizzo di una cache locale a Kong.\\
Il plugin, insieme al Gateway, funge quindi da tramite tra l'utente e il microservizio, intercettando tutte le richieste dell'utente, analizzandole, scomponendole e 
inoltrandole al microservizio, che risponderà al plugin che a sua volta inoltrerà la risposta all'utente.\\ \\
Quindi, il progetto ha tre macro – componenti:
\begin{itemize}
\item Kong Gateway;
\item Plugin per Kong Gateway scritto in linguaggio Lua;
\item Microservizio scritto in linguaggio Java.
\end{itemize}

L'obiettivo del progetto è quello di realizzare un software che consenta l'autenticazione di un utente all'interno di una piattaforma generica tramite un token 
da includere all'interno della richiesta inviata, senza quindi l'utilizzo di password o altri strumenti di autenticazione, 
come dettagliato nel \emph{Capitolo \ref{chap:requisiti} – Requisiti e casi d'uso}. \\

Si è deciso, insieme al Tutor Aziendale, di puntare su Kong Gateway come API Gateway per questo progetto proprio perché consente l'installazione di plugin 
proprietari e custom, oltre a fornire tutte le funzioni di base di un API Gateway.\\
Per quanto riguarda lo sviluppo del plugin, si è voluto utilizzare il linguaggio Lua per le sue caratteristiche di semplicità di utilizzo e integrazione e leggerezza.\\ \\

Per la realizzazione del microservizio è stato deciso di utilizzare il linguaggio Java e il framework Spring perché sono stati reputati più adatti allo scopo.\\
Il framework Spring fornisce delle caratteristiche che risultano essere comode e sicure per il programmatore per quanto riguarda la creazione di microservizi, 
gestione dei dati, integrazione con un Database ecc., il tutto mantenendo intatti i costrutti e le pratiche del linguaggio Java.\\
Tutte le tecnologie utilizzate all'interno del progetto sono approfondite nel \emph{Capitolo \ref{chap:tecnologie} – Tecnologie utilizzate}.\\ \\

Più nello specifico, la realizzazione del progetto si è articolata in diverse fasi, iniziando dallo sviluppo del Database e del microservizio in Java.\\
Lo sviluppo è stato effettuato all'inizio in locale, con i relativi test di corretto funzionamento e, successivamente, si è passati all'installazione di Kong Gateway
sulla Virtual Machine creata su Microsoft Azure e quindi allo sviluppo del plugin in Lua per il Gateway.\\
Infine, si è posto il microservizio all'interno della Virtual Machine (con il relativo ambiente di esecuzione, la JVM) e quindi sono state finalizzate tutte le 
impostazioni per la corretta comunicazione tra microservizio, plugin, Gateway e utente.\\
Per la fase di test, sia in locale che in remoto, è stato utilizzato il software Postman.\\
Tutti questi argomenti, lo sviluppo e il testing, saranno approfonditi rispettivamente nel \emph{Capitolo \ref{chap:progettazionesviluppo} – Progettazione e sviluppo} 
e \emph{Capitolo \ref{chap:testing} – Testing}.


%
%%%% I Capitoli di Contenuto	
\pagestyle{fancy}

\chapter{Descrizione dello stage}\label{chapter:formattazione}
\section{Progetto di stage}\label{sec:cap_sec_subsec}
\section{Azienda}\label{sec:images}
\section{Introduzione al progetto}\label{sec:code}
\section{Obiettivi}
% Prima di introdurre il capitolo si può scrivere una breve introduzione su ciò che si andrà ad affrontare.
% In questo capitolo, per esempio, saranno presentati alcuni esempi di formattazione degli oggetti più comuni di una Tesi in informatica.


% \section{Progetto di stage}\label{sec:cap_sec_subsec}
% Capitoli, sezioni e sottosezioni devono essere usate appropriatamente e non sostituire elenchi puntanti.
% In particolare, i Capitoli devono riguardare macro-argomenti della tesi; ad esempio \textbf{Background}, \textbf{Obiettivi}, \textbf{Progettazione/Implementazione} e \textbf{Risultati}.
% Questi capitoli sono semplicemente una traccia bisogna adattare a seconda delle esigenze.

% Le sezioni invece devono riguardare argomenti all'interno della macro-area definita dal capitolo.
% Ad esempio se più tecniche sono state utilizzate durante la tesi si può suddividere il capitolo \textbf{Progettazione/Implementazione} in sezioni, ognuna riguardante una delle tecniche sperimentate.

% Le sottosezioni infine si possono utilizzare per descrivere concetti distinti all'interno di ogni sezione, ognuno dei quali deve avere una sua identità.
% Ogni sottosezione deve avere un motivo per essere definita e non semplicemente separare parti di uno stesso discorso, per quello c'è l'indentazione (doppio ``a capo'' per separare parti distinte dello stesso paragrafo e ``$\backslash\backslash$'' per separare due paragrafi).

% %\subsection{Esempio di sottosezione}\label{subsec:es_subsec}
% %\lipsum[1-2]

% \section{Azienda}\label{sec:images}
% Di seguito alcuni esempi di immagini.
% Figura~\ref{fig:one} presenta una singola immagine con ancoraggio ``\textbf{ht}" che chiede al altex di lasciare la figura dove si trova, se possibile, e altrimenti di metterla in alto alla pagina.
% Figura~\ref{fig:two} presenta una singola immagine con due sotto-figure con ancoraggio ``\textbf{ht}".
% Figura~\ref{fig:three} presenta una singola immagine con tre sotto-figure con ancoraggio ``\textbf{H}" che forza l'immagine nel posto scelto dall'utente, questa opzione è sconsigliata a meno di necessità particolari.
% Ogni volta che un'immagine è inserita è buona norma citarla nel testo, altrimenti l'immagine non avrebbe un chiaro significato.
% Nel caso delle sotto-immagini si possono citare anche loro, ad esempio Figura~\ref{fig:lev1} è composta da altre due sotto-figure.

% \section{Introduzione al progetto}\label{sec:code}
% Alcuni esempi di codice, definizioni, tabelle e altro.

% \begin{algorithm}[ht]
% 	\caption{Esempio di pseudo-codice}
% 	\label{alg:Prim_Mst}
% 	\begin{algorithmic}[1]
% 		\Statex
% 		\Function{MST-Prim}{grafo G, funzione\_peso $\omega$, nodo\_radice r}\\
% 		\State{\textit{Q}: coda di priorita contenente tutti i vertici in \textit{V}}
% 		\For{ogni \textit{u} $\in$ V(G)}
% 		\Let{\textit{u.key}}{$\infty$}
% 		\Let{\textit{u.}$\pi$}{\textit{NIL}} 
% 		\EndFor
% 		\Let{\textit{r.key}}{0}
% 		\Let{\textit{Q}}{\textit{V(G)}}
% 		\While{\textit{Q} $\neq$ 0}
% 		\Let{\textit{u}}{EXTRACT-MIN(\textit{Q})}
% 		\For{ogni \textit{v} $\in$ \textit{G.Adj[u]}}
% 		\If{$v \in Q$ and $\omega(u,v) < v.key$}
% 		\Let{\textit{v.key}}{$\omega(u,v)$}
% 		\Let{\textit{v.}$\pi$}{\textit{v}}
% 		\EndIf
% 		\EndFor
% 		\EndWhile
% 		\State \Return{}
% 		\EndFunction
% 	\end{algorithmic}
% \end{algorithm}


% \begin{lstlisting}[caption={Esempio di definizione struttura dati C++, ma non algoritmo},
% 	label={lst:block_struct}]
% 	struct Block
% 	{
% 		int id_block;
% 		char resolution;
% 		int id_subdomain;
% 		int key;
% 		bool in_other_subdomain;
% 		struct neigh_t neighbors[4];
% 	};
% \end{lstlisting}



% \begin{algorithm}[ht]
% 	\caption{Esempio di codice imperativo}
% 	\label{alg:my_prim_multi}
% 	\begin{algorithmic}[1]
% 		\Statex
% 		\Function{PRIM\_MULTI-MST}{roots\_list \textit{roots}, array\_Block \textit{adjacency\_list}}
		
% 		\State{\color{blue}{//Inizializzazione di tutti i vertici}}
% 		\For{ogni \textit{u} $\in$ \textit{adjacency\_list}}
% 		\Let{\textit{u.key}}{$\infty$}
% 		\Let{\textit{u.id\_subdomains}}{\textit{NIL}}
% 		\Let{\textit{u.in\_other\_subdomain}}{FALSE} 
% 		\EndFor
		
% 		\State{\color{blue}{//Inizializzazione di tutte le radici}}
% 		\Let{\textit{i}}{0}
% 		\For{ogni \textit{r} $\in$ \textit{roots}}
% 		\Let{\textit{r.key}}{0}
% 		\Let{\textit{r.id\_subdomains}}{\textit{i}}
% 		\Let{\textit{i}}{$i+1$}
% 		\EndFor
% 		\State{\textit{Q}: coda di priorita ordinata in base al campo 
% 			\textit{key}}
% 		\Let{\textit{Q}}{\textit{roots}}
% 		\Let{\textit{count}}{0}
		
% 		\While{count $<$ \textit{adjacency\_list.size}}
% 		\Let{\textit{u}}{EXTRACT-MIN(\textit{Q})}
		
% 		\If{!(\textit{u.in\_other\_subdomain})}
% 		\State{\color{blue}{//Essendo il minimo viene aggiunto definitivamente}}
% 		\Let{\textit{u.in\_other\_subdomain}}{TRUE}
% 		\State{\color{blue}{//Si analizzano i vicini}}
% 		\For{ogni \textit{v} $\in$ \textit{u.neighbors}}
% 		\State{\color{blue}{//In questo caso specifico il peso di ogni arco vale 1}}
% 		\If{$v \in adjacency\_list$ and ($u.key+1) < v.key$}
% 		\Let{\textit{v.key}}{$u.key+1$}
% 		\Let{\textit{v.subdomains}}{\textit{u.subdomains}}
% 		\Let{\textit{count}}{\textit{count} + 1}
% 		\EndIf
% 		\EndFor
% 		\EndIf
% 		\EndWhile
% 		\State \Return{}
% 		\EndFunction
% 	\end{algorithmic}
% \end{algorithm}

% \begin{algorithm}[ht]
% 		\caption{esempio di algoritmo in C++}
% 		\label{lst:genic_mpi}
% \begin{lstlisting}
% 	#include (*@\textquotedblleft@*)mpi.h"
	
% 	(*@~ ~ ~ ~\raisebox{-1pt}[0pt][0pt]{$\vdots$}@*)
	
% 	main(int argc, char** argv){
		
% 		(*@~ ~ ~ ~ ~ ~ ~ ~ ~ ~{\raisebox{-1pt}[0pt][0pt]{$\vdots$}}@*)
		
% 		//Nessuna chiamata a funzioni MPI prima di questa
% 		MPI_Init(&argc, &argv);
		
% 		(*@~ ~ ~ ~ ~ ~ ~ ~ ~ ~{\raisebox{-1pt}[0pt][0pt]{$\vdots$}}@*)
		
% 		MPI_Finalize();
% 		//Nessuna chiamata a funzioni MPI dopo questa
		
% 		(*@~ ~ ~ ~ ~ ~ ~ ~ ~ ~{\raisebox{-1pt}[0pt][0pt]{$\vdots$}}@*)
% 	}	
% \end{lstlisting}
% \end{algorithm}
	
% \begin{defn}[Esempio di Definizione]
% Un piccolo esempio di definizione
% \end{defn}

% \begin{table}[ht]
% 	\centering
% 	\resizebox{0.99\linewidth}{!}{
% 		\begin{tabular}{||c||c|c|c|c|c|c|c|c|c|c||}
% 			\hhline{|t:=:t:==========:t|}
% 			CT scan &\phantom{00}01\phantom{00}&\phantom{00}02\phantom{00}&\phantom{00}03\phantom{00}&\phantom{00}04\phantom{00}&\phantom{00}05&\phantom{00}06\phantom{00}&\phantom{00}07\phantom{00}&\phantom{00}08\phantom{00}&\phantom{00}09\phantom{00}&\phantom{00}10\phantom{00} \\
% 			\hhline{|:=::==========:|}
% 			Input arteries (n.) &66&175&76&134&198&172&154&108&91&65 \\
% 			\hhline{||-||----------||}
% 			Labeling time (s)&19.13&87.38&26.11&47.08&90.89&356.95&303.77&91.95&16.01&21.88 \\
% 			\hhline{||-||----------||}
% 			Grounding time (s)&9.76&52.45&9.57&22.88&70.34&47.52&37.44&17.98&10.61&8.04 \\
% 			\hhline{||-||----------||}
% 			Optimum time (s)&6.66&31.05&14.67&20.63&15.2&29.17&54.71&61.03&4.28&7.24 \\
% 			\hhline{|b:=:b:==========:b|}
% 		\end{tabular}
% 	}
% 	\vspace*{2mm}
% 	\caption{Esempio di tabella}
% 	\label{tab:perf}
% \end{table}


% \section{Obiettivi}
% Durante la scrittura della Tesi è necessario fornire i riferimenti bibliografici che hanno permesso la stesura della tesi stessa.
% Questo si ottiene semplicemente tramite l'utilizzo del comando ``cite'' che prende come argomento la label di una delle entry nel file .bib.
% Quindi per citare una qualunque fonte, ad esempio ``Artificial Intelligence - A Modern Approach", si può fare \cite{modernApproach} ($\backslash\text{cite}\{\text{modernApproach}\}$).
% Questo leggerà la referenza nel file .bib che ha come label ``cormen".
% Per citare più riferimenti contemporaneamente basta sperare le varie labels con una virgola come segue \cite{gelfond1998action,modernApproach,durfee1999distributed,de2003resource,allen2009complexity,bernstein2002complexity}.
% Le reference non citate non appariranno in bibliografia.
% Google scholar e dblp sono siti nei quali estrarre la reference per bibtex per un lavoro di cui si conosce solo il titolo, ad esempio.
% % !TeX spellcheck = it_IT

\chapter{Tecnologie utilizzate}\label{sec:second_mpi}
\section{Git}
\section{Java}
\section{Eclipse}
\section{Framework Spring}
\subsection{Spring Boot}
\subsection{Spring Data}
\section{Azure}
\section{MySQL}
\section{Kong Gateway}
\section{Lua}
\section{Docker}
\section{Postman}

% \emph{MPI} cioè, Message Passing Interface, è uno \textit{standard per l'interfaccia di una libreria per lo scambio di messaggi} \cite{mpi_standard}.\\
% Questo standard \textit{de facto} è nato grazie al lavoro svolto su \textit{MPI Forum}, un forum nel quale aziende che si occupano di \textquotedblleft Parallel Computing\textquotedblright, ricercatori nell'ambito informatico, sviluppatori di librerie e di applicazioni hanno collaborato per raggiungere questo risultato.\\
% Come detto \textit{MPI} è una specifica e non una libreria e si propone quindi di imporre regole e di dare istruzioni su quello che una libreria per lo scambio di messaggi deve fare. Inoltre, non essendo \textit{MPI} un linguaggio di programmazione tutte le operazioni sono espresse come funzioni, subroutine, o metodi a seconda del \textit{binding} con i linguaggi che fanno parte dello standard.\\
% L'obbiettivo di \textit{MPI}, quindi, è quello di stabilire uno standard per il \textit{message-passing} \emph{portabile}, \emph{flessibile}, \emph{pratico} ed \emph{efficiente}.\\
% \textit{MPI} è rivolto principalmente al modello di \textit{message-passing} nella \textit{programmazione parallela}, in cui i dati vengono spostati dallo spazio di indirizzi di un processo a quello di un altro tramite operazioni cooperative su ciascun processo.\\
% In particolare in un ambiente di comunicazione a memoria distribuita, in cui le astrazioni di livello superiore sono costruite sulla base del livello inferiore, ovvero sulle routine di \textit{message-passing}, i benefici della standardizzazione sono particolarmente evidenti.\\
% Inoltre, la definizione di uno standard \textit{message-passing}, come quella proposta da MPI, fornisce un insieme di funzioni base ben definito, efficientemente implementabili dai vendor interessati al \textquotedblleft Parallel Computing\textquotedblright, o per le quali sia possibile fornire supporto hardware, rafforzando in tal modo la scalabilità.\\
% L'obiettivo di \textit{MPI} è quindi quello di sviluppare uno standard ampiamente utilizzato per la scrittura di programmi di \textit{message-passing}.

% \subsection{La struttura di un programma MPI}
% Un programma \textit{MPI} è un programma che, seguendo le procedure introdotte dallo standard, esegue operazioni di \textit{message-passing} su architetture di calcolo parallelo.\\
% D'ora in avanti si affronteranno programmi \textit{MPI} in \textit{C/C++} che utilizzano la libreria \emph{OpenMPI} \cite{open_mpi} per l'implementazione delle funzioni, ma le considerazioni fatte sono comunque valide per le relative funzioni legate al \textit{Fortran} e per implementazioni di altre librerie.\\
% Ogni programma \textit{MPI} deve utilizzare la direttiva del preprocessore: \lstinline{#include "mpi.h"}\\
% \textit{mpi.h} infatti contiene le definizioni, le macro e i prototipi delle funzioni necessari per compilare un programma \textit{MPI}.\\
% Prima di ogni altra funzione \textit{MPI} bisogna invocare la funzione \textit{MPI\_Init()}; questa funzione deve essere chiamata una sola volta. La \textit{MPI\_Init()} prende come argomenti i puntatori ai parametri del \textit{main} e permette al sistema di preparare l'ambiente per l'esecuzione delle varie routine \textit{MPI} che possono essere usate.\\
% Una volta che il programma ha finito di utilizzare le funzioni \textit{MPI} deve chiamare la \textit{MPI\_Finalize()} che permette di \textquotedblleft pulire\textquotedblright~ogni lavoro non terminato dalle varie funzioni \textit{MPI} invocate precedentemente.\\
% Nel listato \ref{lst:genic_mpi} è mostrata la generica struttura di un programma \textit{MPI} \cite{mpi_general}.\\
% \begin{lstlisting}[label={lst:genic_mpi},
% caption={Programma MPI generico}]

% (*@~ ~ ~ ~\raisebox{-1pt}[0pt][0pt]{$\vdots$}@*)

% #include (*@\textquotedblleft@*)mpi.h"

% (*@~ ~ ~ ~\raisebox{-1pt}[0pt][0pt]{$\vdots$}@*)

% main(int argc, char** argv){

% (*@~ ~ ~ ~ ~ ~ ~ ~ ~ ~{\raisebox{-1pt}[0pt][0pt]{$\vdots$}}@*)

%     //Nessuna chiamata a funzioni MPI prima di questa
%     MPI_Init(&argc, &argv);

% (*@~ ~ ~ ~ ~ ~ ~ ~ ~ ~{\raisebox{-1pt}[0pt][0pt]{$\vdots$}}@*)

%     MPI_Finalize();
%     //Nessuna chiamata a funzioni MPI dopo questa

% (*@~ ~ ~ ~ ~ ~ ~ ~ ~ ~{\raisebox{-1pt}[0pt][0pt]{$\vdots$}}@*)
% }

% \end{lstlisting}
% \subsection{Le principali componenti}\label{subsec:comm_mpi}
% In MPI la comunicazione fra i processi si basa sullo scambio di messaggi.\\
% MPI quindi definisce le strutture e le informazioni che servono per realizzare uno scambio di messaggi tra processi.\\
% In particolare lo standard descrive:
% \begin{itemize}
% 	\item In che modo rappresentare il dato, da inviare nel messaggio, durante la comunicazione;
% 	\item Come specificare i processi protagonisti dello scambio del messaggio;
% 	\item I vari modi in cui può essere implementata la comunicazione.
% \end{itemize}  
% \subsubsection{Definizione dei dati in MPI}\label{subsec:def_dati_mpi}
% Per poter inviare un dato attraverso un messaggio, che poi verrà scambiato tra i processi, occorre definire alcune informazioni fondamentali legate al dato stesso; in particolare è necessario definire:
% \begin{itemize}
% 	\item \textbf{Il buffer}, ovvero l'indirizzo di memoria che contiene i dati da inviare;
% 	\item \textbf{Il numero di elementi} da inviare nel messaggio;
% 	\item \textbf{Il tipo di dato} degli elementi.
% \end{itemize}
% Il tipo di dato può rispecchiare uno dei tipi di dato più comuni del linguaggio \textit{C} (figura \ref{fig:mpi_data_types}) oppure può essere un tipo di dato definito da utente in quanto MPI fornisce le primitive per costruire strutture dati partendo da altre già esistenti.

% \subsubsection{Specificazione dei processi coinvolti nella comunicazione}\label{subsec:comm_rank_tag}
% Durante l'invio di un messaggio è necessario specificare il destinatario e garantire che le informazioni che lo identificano siano univoche per non compromettere la comunicazione.\\
% A questo scopo durante l'invio di un messaggio MPI bisogna sempre indicare:
% \begin{itemize}
% 	\item \textbf{Il communicator} che rappresenta uno specifico contesto a cui appartiene un gruppo di processi (il \textit{communicator} di default si chiama \textit{MPI\_COMM\_WORLD} e racchiude tutti i processi disponibili).
% 	\item \textbf{Il rank}, che è un identificativo univoco assegnato ad ogni processo, del processo destinatario;
% 	\item \textbf{Il tag} che permette, grazie ad un intero definito da utente, di differenziare ulteriormente il messaggio.
% \end{itemize}

% \subsubsection{Le implementazioni delle comunicazioni punto a punto}
% MPI prevede diverse modalità per gestire la comunicazione tra i processi.
% \paragraph{La comunicazione standard}~\newline~\newline
% La modalità standard è \emph{bloccante e asincrona} si effettua utilizzando la funzioni \textit{MPI\_Send} per l'invio e \textit{MPI\_Recv} per la ricezione.\\
% \begin{lstlisting}[label={lst:mpi_send_recv},
% caption={Intestazioni delle primitive di comunicazione MPI}]
% MPI_Send(void *buf, int count, MPI_Datatype datatype, int dest, int tag, MPI_Comm comm)
         
% MPI_Recv(void *buf, int count, MPI_Datatype datatype, int source, int tag, MPI_Comm comm, MPI_Status *status)
% \end{lstlisting}

% In entrambe le funzioni i parametri \textit{buf}, \textit{count} e \textit{datatype} servono per fornire le informazioni sul contenuto nel messaggio come spiegato nel paragrafo \ref{subsec:def_dati_mpi}.\\
% Il parametro \textit{dest} nella funzione di invio e il parametro \textit{source} in quella di ricezione indicano rispettivamente il \textit{rank} del processo che deve ricevere e di quello che deve inviare i dati. Oltre al \textit{rank} dei processi coinvolti, queste funzioni prendono in argomento anche le altre informazioni illustrate nel paragrafo \ref{subsec:comm_rank_tag} relative al destinatario; in particolare richiedono: il \textit{tag} del messaggio e il \textit{communicator(comm)} nel quale si svolge la comunicazione.\\
% La funzione \textit{MPI\_Recv} grazie a \textit{status} può ottenere informazioni aggiuntive sul messaggio.\\
% Essendo la \textit{MPI\_Send} bloccante, ritorna solo quando i dati sono stati copiati nel buffer di sistema del destinatario e quindi, per esempio, se il buffer è pieno, la funzione resta bloccata in attesa che si liberi. Una volta che la funzione è ritornata, non c’è alcuna garanzia che il processo destinatario abbia ricevuto il messaggio.\\
% A sua volta la \textit{MPI\_Recv}, essendo bloccante ritorna solamente quando i dati sono stati copiati dal buffer di sistema all’indirizzo di memoria specificato.\\
% Lo schema di una comunicazione bloccante è mostrato in figura \ref{fig:mpi_block_com}.
% \paragraph{Altri tipi di comunicazione punto a punto}~\newline~\newline
% Oltre alla modalità standard (asincrona) MPI definisce altre modalità, \textit{bloccanti e non}, di comunicazione per l’invio di un messaggio.\\
% Le seguenti modalità di invio sono \emph{bloccanti} e richiedono come parametri gli stessi richiesti dalla funzione primitiva MPI\_Send:
% \begin{itemize}
% 	\item MPI\_Ssend ($\cdots$): questa funzione ritorna solo quando il destinatario ha iniziato a ricevere;
% 	\item MPI\_Bsend ($\cdots$): questa funzione ritorna quando i dati sono copiati in un buffer specifico;
% 	\item MPI\_Rsend ($\cdots$): da utilizzare solo quando si è certi che il destinatario sia in ascolto.
% \end{itemize}
% Esistono anche modalità di comunicazione \emph{non bloccanti} (di cui viene mostrato lo schema in figura \ref{fig:mpi_non_block_comm}) che vengono descritte dalle seguenti funzioni nel codice \ref{lst:mpi_send_recv_non_blocc}.\\
% \begin{lstlisting}[ 
% label={lst:mpi_send_recv_non_blocc},
% caption={Intestazioni delle comunicazioni non bloccanti MPI}]
% MPI_Isend(void* buf, int count, MPI_Datatype datatype, int dest, int tag, MPI_Comm comm, MPI_Request &request)

% MPI_Irecv(void* buf, int count, MPI_Datatype datatype, int source, int tag, MPI_Comm comm, MPI_Request &request)
% \end{lstlisting}
% Le funzioni sopracitate ritornano immediatamente, ma la memoria all’indirizzo \textit{buf} non può essere scritta o letta finché non si è certi della conclusione dell’operazione.\\
% Il campo \textit{request} delle due funzioni è un \textit{handle} che permette di verificare e interagire con lo status della richiesta grazie ad opportune funzioni (MPI\_Test e MPI\_Wait).

% \subsubsection{Le comunicazioni collettive}\label{subsubsec:coll_comm}
% Un ultimo tipo di comunicazione messo a disposizione da MPI è la comunicazione collettiva. Una comunicazione collettiva è una comunicazione che riguarda tutti i processi presenti all'interno di uno stesso \textit{communicator} che non interferisce con le comunicazioni \emph{punto-punto} e che non può essere bloccante.\\
% MPI mette a disposizione diverse funzioni per gestire vari tipi di comunicazione collettiva (figura \ref{fig:mpi_collect_comm}):
% \begin{itemize}
% \item \textbf{MPI\_Bcast} per permettere ad un processo di inviare lo stesso messaggio a tutti gli altri;
% \item \textbf{MPI\_Reduce} per permettere ad un processo di suddividere un messaggio tra tutti gli altri;
% \item \textbf{MPI\_Scatter} per far si che tutti i processi inviino un messaggio ad uno stesso nodo;
% \item \textbf{MPI\_Gather} per far si che tutti i processi cooperino per calcolare un risultato da inviare ad uno stesso processo.
% \end{itemize}


% \subsection{Il modello hardware per un ambiente MPI}\label{subsubsec:mpi_hw_model}
% Come detto MPI è un'interfaccia per le librerie che si devono occupare di message-passing.\\
% Questo scambio di informazioni è dunque legato ad un ambiente di calcolo parallelo e a una relativa architettura hardware che consente di eseguire uno stesso programma parallelamente su più processori o addirittura su più nodi.\\
% Le varie configurazioni che può avere un ambiente di calcolo parallelo, secondo la tassonomia di Flynn (1986), derivano dalla configurazione hardware e software di due componenti principali del sistema: il \emph{flusso di istruzioni (processi)} e il \emph{flusso di dati}. Le varie configurazioni previste da questa suddivisione sono:
% \begin{itemize}
% 	\item \textbf{SISD} (\textit{Single Instruction Single Data}): Architetture dei sistemi seriali, senza quindi, nessun grado di parallelismo (schema in figura \ref{fig:sisd});
% 	\item \textbf{SIMD} (\textit{Single Instruction Multiple Data}): Architetture composte da molte unità di elaborazione che eseguono contemporaneamente la stessa istruzione ma lavorano su insiemi di dati diversi (schema in figura \ref{fig:simd});
% 	\item \textbf{MISD} (\textit{Multiple Instruction Single Data}): Architetture in cui più flussi di istruzioni lavorano contemporaneamente su un unico flusso di dati (schema in figura \ref{fig:misd}). Questa tipologia di architettura non è stata praticamente mai adottata;
% 	\item \textbf{MIMD} (\textit{Multiple Instruction Multiple Data}): Architetture in cui più istruzioni vengono eseguite contemporaneamente su più dati diversi (schema in figura \ref{fig:mimd}). Sotto questa classificazione ricadono i \textit{cluster} di computer e di conseguenza l'ambiente sul quale è stato basato lo sviluppo di questa tesi.
% \end{itemize}

% In particolare tra i sistemi MIMD possono distinguersi le architetture a \emph{memoria condivisa} (figura \ref{fig:mem_dist}) o quelle a \emph{memoria distribuita} (figura \ref{fig:mem_cond}).\\
% Nella prima configurazione la memoria è condivisa fisicamente (UMA) o solo logicamente assegnando indirizzi globali alle vari memorie locali (NUMA), permettendo ai vari nodi di lavorare sugli stessi spazi di memoria, evitando quindi di dover gestire le comunicazioni.\\
% Nell'architettura a memoria distribuita, ogni nodo possiede una propria memoria locale che non fa parte dello spazio di indirizzamento degli altri processori e, per regolare l'accesso di un nodo alle risorse che non sono locali, bisogna utilizzare protocolli di comunicazione e scambio dati come, per esempio, quelli descritti da MPI.
\chapter{Analisi dei requisiti}\label{chapter:formattazione}
\section{Casi d'uso}\label{sec:cap_sec_subsec}
% % !TeX spellcheck = it_IT
% \section{I Grafi e il loro Partizionamento}
% Nel capitolo \ref{sec:first_swe} di questa tesi si è discusso di come le informazioni del territorio, su cui si vuole svolgere la simulazione, siano completamente mappate su strutture dati specifiche del progetto (celle e blocchi).\\
% L'unione dei blocchi (composti a loro volta dalle celle), crea il dominio logico sul quale verranno svolte le operazioni di \emph{partizionamento} spiegate nel capitolo \ref{sec:fourth_sudd_dom}.\\
% Prima di descrivere le operazioni svolte sul dominio è necessario introdurre il concetto di \emph{grafo} e alcune delle operazioni legate ad esso. Il grafo è la struttura di base su cui si basa la rappresentazione del dominio e la relativa suddivisione.

% \subsection{Introduzione ai grafi}
% Il \emph{grafo} è una struttura matematica che ha trovato ampio utilizzo nell'ambito dell'informatica.\\
% Concettualmente, un grafo, è formato da \emph{vertici} e da \emph{archi}, che connettono una coppia di vertici tra loro \cite{cormen}.
% \begin{defn}
% Si dice grafo \emph{G} una coppia \emph{(V,E)} dove dove\emph{V(G)} è l'insieme dei \emph{vertici} (o nodi) ed \emph{E(G)} è l'insieme degli \emph{archi} (che può anche essere vuoto).
% \end{defn}
% Un grafo semplice \emph{G} consiste quindi in un insieme finito e non-vuoto \textit{V(G)} di vertici e in un insieme finito \textit{E(G)} di coppie distinte e non ordinate di elementi di\textit{V(G)}.
% \begin{figure}[H]
% 	\centering
% 	\includegraphics{img/logo_unipr}
% 	\caption{Un semplice esempio di grafo}
% 	\label{fig:simple_graph}
% \end{figure}

% \subsubsection{I vari tipi di grafo}\label{subsubsec:graph_type}
% Quando si parla di grafi bisogna però distinguere di quale tipologia si sta trattando.\\
% Esistono infatti vari tipi di grafo ognuno con caratteristiche differenti \cite{graph_theory}. La prima distinzione si ha tra grafi \emph{diretti} e \emph{indiretti}:

% \paragraph{Grafi Indiretti}~\\
% Nei grafi indiretti la relazione di adiacenza è simmetrica: l'arco che connette il nodo \textit{u} al nodo \textit{v} è lo stesso che connette il \textit{v} al nodo \textit{u}. Pertanto l'ordine dei vertici nella coppia che compone l'arco non ha importanza (figura \ref{fig:udir_graph}).

% \paragraph{Grafi Diretti}~\\
% In un grafo diretto (o orientato) gli elementi di \textit{E(G)} sono coppie \emph{ordinate} e di conseguenza, gli archi sono associati ad una direzione (figura \ref{fig:dir_graph}). Per esempio l'arco dal vertice \textit{u} al nodo \textit{v}, indicato come (u, v), è diverso dall'arco di direzione opposta (v, u).\\~\\

% \begin{defn}
% 	Si definisce \emph{ciclo} un cammino lungo archi orientati che parte da un vertice ed arriva allo stesso.
% \end{defn}

% Un'altra caratterizzazione dei grafi tiene conto della presenza o meno di un ciclo: si distingue tra \textbf{Grafi diretti Ciclici} (figura \ref{fig:dir_graph}) nel caso in cui il ciclo sia presente e \textbf{Grafi diretti Aciclici (DAG)} (figura \ref{fig:dag}) qualora non vi sia.\\

% Infine un'ultima caratteristica che possono avere tutti i tipi di grafo sopra elencati, è l'associazione o meno di un valore, detto \emph{peso}, ad ogni arco appartenete al grafo.\\
% Un \textbf{Grafo Pesato} (figura \ref{fig:weighted_graph}) \textit{G} è quindi un grafo dove ogni elemento dell'insieme \textit{E(G)} ha associato un valore numerico. Questa associazione permette di eseguire operazioni sui vari archi tramite il confronto dei pesi, come ad esempio, individuare un cammino minimo tra due vertici.

% \begin{figure}[H]
% 	\centering
% 	\begin{subfigure}{0.3\textwidth}
% 		\centering
% 		\includegraphics[width=0.8\linewidth]{img/logo_unipr}
% 		\caption{Esempio di grafo indiretto}
% 		\label{fig:udir_graph}
% 	\end{subfigure}%
% \begin{subfigure}{0.3\textwidth}
% 	\centering
% 	\includegraphics[width=0.8\linewidth]{img/logo_unipr}
% 	\caption{Esempio di grafo diretto che contiene un ciclo}
% 	\label{fig:dir_graph}
% \end{subfigure}%
% 	\begin{subfigure}{0.4\textwidth}
% 		\centering
% 		\includegraphics[width=0.8\linewidth]{img/logo_unipr}
% 		\caption{Esempio di DAG (grafo diretto aciclico)}
% 		\label{fig:dag}
% 	\end{subfigure}
% 	\begin{subfigure}{1.0\textwidth}
% 		\centering
% 		\begin{subfigure}{0.5\textwidth}
% 			\centering
% 			\includegraphics[width=0.7\linewidth]{img/logo_unipr}
% 		\end{subfigure}%
% 		\begin{subfigure}{0.5\textwidth}
% 			\centering
% 			\includegraphics[width=0.7\linewidth]{img/logo_unipr}
% 		\end{subfigure}
% 		\caption{Due esempi di grafi pesati}
% 		\label{fig:weighted_graph}
% 	\end{subfigure}
% 	\label{fig:graph_flavours}
% 	\caption{I diversi tipi di grafo}
% \end{figure}


% \subsection{La rappresentazione dei grafi}\label{subsec:graph_rappr}
% I grafi sono utilizzati spesso per descrivere entità fisiche all'interno della memoria di un computer. Nel caso di questa tesi viene utilizzato un grafo orientato per descrivere il territorio nel quale si svolge la simulazione.\\
% Esistono due metodi standard per rappresentare un grafo $G = (V, E)$: attraverso una \emph{collezione di liste di adiacenza} e attraverso una \emph{matrice di adiacenza}.

% \subsubsection{Rappresentazione con Matrice di adiacenza}
% Per questa rappresentazione si suppone che i vertici del grafo $G = (V, E)$ siano numerati da 1 a $\arrowvert V\arrowvert$ in modo arbitrario. La rappresentazione tramite matrice di adiacenza del grafo consiste in una matrice di dimensioni $\arrowvert V\arrowvert \times \arrowvert V\arrowvert$ in cui vale che l'elemento $m_{ij}$ rispetti la seguente condizione:\\
% \begin{equation*}
% m_{ij} = \begin{cases}
% 1 &\text{se ($i,j$) $\in E$}\\
% 0 &\text{altrimenti}
% \end{cases}
% \end{equation*}
% La matrice di adiacenza, mostrata in figura \ref{fig:matrix_adj} occupa una quantità di memoria quadratica rispetto al numero di vertici indipendentemente dal numero degli archi del grafo.\\
% La rappresentazione di un grafo pesato tramite matrice di adiacenza, sostituisce il valore 1 che indica presenza dell'arco tra i nodi con il peso dell'arco stesso e utilizza un valore speciale (come 0, $\infty$ o \textit{NIL}) per il caso in cui l'arco non sia presente.\\
% La matrice di adiacenza è consigliata quando il grafo è un \emph{grafo denso}  ovvero $\arrowvert E\arrowvert$ si avvicina a $\arrowvert V\arrowvert ^2$ o quando c'è bisogno di constatare rapidamente se esiste un arco specifico.

% \subsubsection{Rappresentazione con Liste di adiacenza}
% La rappresentazione tramite liste di adiacenza di un grafo $G = (V, E)$ consiste in  un array (indicato d'ora in avanti con \textit{Adj}) di $\arrowvert V\arrowvert$ liste, una per ogni vertice. Per ogni nodo quindi si crea una lista che contiene tutti i vertici adiacenti. Un esempio di lista di adiacenza è mostrato in figura \ref{fig:list_adj}.\\
% La memoria richiesta da questa rappresentazione dipende dal tipo di grafo su cui si basa: infatti la somma totale delle lunghezze delle liste è pari a $\arrowvert E\arrowvert$ se il grafo è orientato, e 2$\arrowvert E\arrowvert$ se il grafo non lo è. La rappresentazione in lista di adiacenza conserva comunque, in entrambi i casi, l'interessante proprietà di occupare una quantità di memoria massima $\Theta(E)$.\\
% Questa rappresentazione è consigliata per la maggioranza dei grafi perchè permette di rappresentare in modo compatto i \emph{grafi sparsi}, quelli in cui $\arrowvert E\arrowvert$ è molto più piccolo di $\arrowvert V\arrowvert ^2$. Inoltre, tramite le liste di adiacenza, è possibile associare facilmente ad ogni arco un peso nel caso in cui si voglia rappresentare un grafo pesato.\\
% L'unico svantaggio è rappresentato dall'impossibilità di determinare velocemente la presenza di uno specificato arco, cosa che invece è possibile fare nelle matrici da adiacenza al costo però di un utilizzo maggiore di memoria.
% \begin{figure}[H]
% 	\centering
% 	\begin{subfigure}{1.0\textwidth}
% 		\centering
% 		\includegraphics[width=0.3\linewidth]{img/logo_unipr}
% 		\caption{Il grafo da rappresentare}
% 	\end{subfigure}
% 	\begin{subfigure}{1.0\textwidth}
% 		\centering
% 		\begin{subfigure}{0.5\textwidth}
% 			\centering
% 			\includegraphics[width=0.7\linewidth]{img/logo_unipr}
% 			\caption{Rappresentazione tramite matrice di adiacenza}
% 			\label{fig:matrix_adj}
% 		\end{subfigure}%
% 		\begin{subfigure}{0.5\textwidth}
% 			\centering
% 			\includegraphics[width=0.7\linewidth]{img/logo_unipr}
% 			\caption{Rappresentazione tramite lista di adiacenza}
% 			\label{fig:list_adj}
% 		\end{subfigure}
% 	\end{subfigure}
% 	\label{fig:adj_matrix_graph}
% 	\caption{Rappresentazione di uno stesso grafo tramite lista e matrice di adiacenza}
% \end{figure}

% \subsection{Il Partizionamento}\label{subsec:partitioning}
% Il \textit{Graph-Partitioning} è un processo di pre-computazione fondamentale per la maggior parte delle applicazioni su larga scala che sono implementate da piattaforme di calcolo parallelo. Il problema del partizionamento definito su un grafo \textit{G} sta nel suddividere \textit{G} in componenti più piccole con specifiche proprietà. Per esempio, una partizione a $k$ divide l'insieme dei vertici in \textit{k} componenti più piccole.\\
% Una partizione è considerata buona se il numero di archi che intercorre tra le componenti separate è piccolo.\\
% Costruire \emph{un grafo a partizione uniforme} è un problema di partizionamento che consiste appunto nel dividere il grafo in componenti della stessa grandezza (in termini di nodi o di peso di questi) cercando di minimizzare le comunicazioni tra le varie componenti.

% %\subsubsection{La complessità}
% %Tipicamente, il problema di partizionamento uniforme ricade sotto la categoria dei problemi \textit{NP-hard}. Le soluzioni a questi problemi sono spesso ricavate mediante l'utilizzo di algoritmi euristici e di approssimazione. Nonostante ciò, il problema di partizionamento uniforme o bilanciata può essere ricondotto ad un problema \textit{NP-completo}.

% \subsection{Alberi di copertura Minimi}\label{subsec:mst_description}
% \begin{defn}
% Dato un grafo $G = (V, E)$ non orientato e connesso, un \emph{albero di copertura} di G è un sottoinsieme $T \subseteq E$ tale che il sottografo S = (V, T) è un albero non orientato connesso ed aciclico.
% \end{defn}
% Si osserva quindi che un albero di copertura è un sottografo di G che contiene tutti i suoi vertici ed è sia aciclico che connesso.\\
% \begin{equation}
% \omega(T) = \sum_{e \text{ } \in \text{ } T} \omega(e)
% \end{equation}
% \begin{defn}
% 	Dato un grafo $G = (V, E)$ non orientato, connesso e pesato sugli archi, un \emph{albero di copertura minimo} di G è un albero di copertura di G di costo minimo. Si abbrevierà il nome in \emph{MST}, dall’inglese Minimum Spanning Tree.
% \end{defn}
% \begin{figure}[H]
% 	\centering
% 	\includegraphics[width=1.0\textwidth]{img/logo_unipr}
% 	\caption{Esempio di un albero di copertura minimo (MST)}
% 	\label{fig:mst_simple}
% \end{figure}

% \subsubsection{Il teorema Fondamentale degli alberi a copertura minimi \cite{cormen}}\label{subsubsec:th_mst}
% \begin{defn}
% Un \emph{taglio} $(S,V\backslash S)$ di un grafo non orientato $G = (V, E)$ è una partizione di $V$.
% \end{defn}
% \begin{defn}
% Un \emph{arco attraversa} il taglio $(S,V\backslash S)$ se uno dei suoi estremi si trova in $S$ e l’altro in $V\backslash S$.
% \end{defn}
% \begin{defn}
% 	Un \emph{taglio rispetta} un insieme $A$ di archi se nessun arco di $A$ attraversa il taglio.
% \end{defn}
% \begin{defn}
% 	Un arco che attraversa un taglio si dice \emph{leggero} se ha peso pari al minimo tra i pesi di tutti gli archi che attraversano tale taglio.
% \end{defn}
% \begin{defn}
% 	Sia $A$ sottoinsieme di un qualche albero di copertura minimo; un arco si dice \emph{sicuro} per $A$ se può essere aggiunto ad $A$ e quest'ultimo continua ad essere sottoinsieme di un albero di copertura minimo.
% \end{defn}
% Siccome le strategie per la ricerca di un albero a connessione minima sono strategie \textit{greedy}, è importante avere un metodo per verificare che gli archi selezionati passo a passo siano \emph{archi validi}. Il teorema di seguito enunciato permette di verificare questa richiesta.
% \begin{thm}(Fondamentale). Sia G = (V, E, $\omega$) un grafo non orientato, connesso e pesato; sia inoltre:
% 	\begin{itemize}
% 	\item $A \subseteq E$ contenuto in qualche albero di copertura minimo;
% 	\item $(S,V\backslash S)$ un taglio che rispetta $A$;
%     \item $(u, v) \in E$ un arco leggero che attraversa il taglio. 
% 	\end{itemize}
% Allora $(u, v)$ è sicuro per $A$.
% \end{thm}

% \subsubsection{L'algoritmo di Prim}
% Introdotto il teorema fondamentale (paragrafo \ref{subsubsec:th_mst}) si può passare a descrivere velocemente l'\emph{algoritmo di Prim} che è alla base delle operazioni di partizionamento (capitolo \ref{sec:fourth_sudd_dom}) svolte in questa tesi.\\
% L’algoritmo di Prim si basa sul teorema fondamentale degli alberi di copertura minimi (paragrafo \ref{subsubsec:th_mst}). Il costo dell'algoritmo varia a seconda della struttura dati di supporto utilizzata: se realizzato con coda di priorità ha un costo di $O(E\cdot log(V))$ mentre se realizzato tramite \textit{heap di Fibonacci} ha un costo di esecuzione di $O(E+V\cdot log(V))$.\\
% Infatti in questo algoritmo gli archi dell'insieme $A$ (ovvero degli archi già selezionati dall'algoritmo) formano sempre un albero singolo. L’algoritmo parte da un nodo scelto $r \in V$, che sarà la radice dell’albero di copertura, e aggiunge ad $A$ un \textit{arco leggero} che collega $A$ con un \textit{vertice isolato} (che non sia estremo di qualche altro arco in A) ad ogni passo, garantendo così che ogni arco aggiunto sia \textit{sicuro}.\\
% Per gestire l’insieme dei vertici non ancora inclusi nell’albero di copertura, viene utilizzata una struttura dati del tipo \textit{coda a min-priorità} basata sul campo \textit{key}.\\
% Il campo \textit{key} è dato dal peso minimo di un arco che collega un vertice \textit{v} a qualsiasi altro nell'albero; nel caso in cui questo arco non esista nell'insieme $A$, il valore di \textit{key} sarà $\infty$.\\
% Inoltre per ogni nodo viene indicato il predecessore tramite il campo $\pi$. Grazie proprio ai predecessori è possibile ricostruire l'albero finale in $A$.
% \begin{equation}
% 	A = \{(v, v.\pi) : v \in (V \backslash \{r\})\}
% \end{equation}

% \begin{algorithm}[H]
% 	\caption{Prim minimun spanning tree}
% 	\label{alg:Prim_Mst}
% 	\begin{algorithmic}[1]
% 		\Statex
% 		\Function{MST-Prim}{grafo G, funzione\_peso $\omega$, nodo\_radice r}\\
% 		\State{\textit{Q}: coda di priorita contenente tutti i vertici in \textit{V}}
% 		\For{ogni \textit{u} $\in$ V(G)}
% 		\Let{\textit{u.key}}{$\infty$}
% 		\Let{\textit{u.}$\pi$}{\textit{NIL}} 
% 		\EndFor
% 		\Let{\textit{r.key}}{0}
% 		\Let{\textit{Q}}{\textit{V(G)}}
% 		\While{\textit{Q} $\neq$ 0}
% 		\Let{\textit{u}}{EXTRACT-MIN(\textit{Q})}
% 		\For{ogni \textit{v} $\in$ \textit{G.Adj[u]}}
% 		\If{$v \in Q$ and $\omega(u,v) < v.key$}
% 		\Let{\textit{v.key}}{$\omega(u,v)$}
% 		\Let{\textit{v.}$\pi$}{\textit{v}}
% 		\EndIf
% 		\EndFor
% 		\EndWhile
% 		\State \Return{}
% 		\EndFunction
% 	\end{algorithmic}
% \end{algorithm}

% \begin{figure}[H]
% 	\centering
% 	\includegraphics[width=0.7\textwidth]{img/logo_unipr}
% 	\caption{Esempio del funzionamento dell'algoritmo di Prim}
% 	\label{fig:prim_alg}
% \end{figure}


\chapter{Progettazione e sviluppo}\label{chapter:formattazione}
\section{Introduzione ai microservizi}\label{sec:cap_sec_subsec}
\section{Introduzione a Kong Gateway}\label{sec:cap_sec_subsec}
\section{Architettura del progetto}\label{sec:cap_sec_subsec}
\section{Architettura del microservizio}\label{sec:cap_sec_subsec}
\subsection{Sviluppo del microservizio}\label{sec:cap_sec_subsec}
\section{Architettura del plugin}\label{sec:cap_sec_subsec}
\subsection{Sviluppo del plugin}\label{sec:cap_sec_subsec}
% % !TeX spellcheck = it_IT

% \section{La ripartizione del Dominio}\label{sec:fourth_sudd_dom}
% Come detto nel paragrafo \ref{subsec:partitioning}, la suddivisione del dominio (che avviene essenzialmente tramite \textit{graph-partitioning}) è un processo di pre-computazione fondamentale per la maggior parte delle applicazioni che lavorano su architetture distribuite. Infatti, un partizionamento corretto permette di sfruttare al meglio la potenza di calcolo offerta da un sistema di \textit{supercomputing}.\\
% Vista la necessità di portare il progetto \textit{Shallow Water Equation on GPU} su un'architettura di nodi distribuita, così da poter raggiungere la potenza di calcolo richiesta per eseguire simulazioni in tempi ragionevoli, si è iniziato a modellare l'applicativo per poterlo adattare alla suddetta architettura.\\
% In particolare questa tesi studia la ripartizione del dominio sui vari nodi.\\
% In questo capitolo verrà quindi introdotto il \textit{criterio di suddivisione} studiato spiegando nel dettaglio la sua implementazione.\\
% La suddivisione verrà affrontata in varie versioni che varieranno i risultati della ripartizione stessa dal punto di vista del bilanciamento (numero di elementi contenuti nei vari sotto-domini) e delle comunicazioni tra i vari sotto-domini.\\
% Si procederà infine a fornire un confronto tra le varie versioni della suddivisione.

% \subsection{La conformazione del dominio}
% Prima di descrivere le varie operazioni di ripartizione è necessario però analizzare nel dettaglio il dominio stesso.\\
% Come detto nel capitolo \ref{sec:first_swe}, il dominio reale del programma per la simulazione numerica dei flussi d’acqua, è una regione di territorio e tutte le informazioni che lo descrivono, vengono mappate nel dominio logico grazie alle celle e ai blocchi.\\
% Ai fini della tesi verranno considerati solo i blocchi, senza tenere conto delle celle che li compongono, in quanto, le informazioni interessanti per la simulazioni vengono fornite dai blocchi stessi e non dalle singole celle che possono pertanto essere tralasciate.\\
% In quest'ottica quando si parla di multi risoluzione ci si riferisce alla differenza di risoluzione dei blocchi e non delle celle che li compongono.\\

% \subsubsection{Le caratteristiche del grafo}
% Per rappresentare la conformazione del dominio reale all'interno del programma si è utilizzato un grafo diretto (paragrafo \ref{subsubsec:graph_type}) in cui l'insieme \textit{V} dei nodi è uguale all'insieme dei blocchi e l'insieme \textit{E} degli archi è composto dalle comunicazioni che ogni blocco ha con i relativi vicini.\\
% La particolare conformazione di questo grafo fa sì che ogni blocco abbia un arco uscente e uno entrante per ogni vicino. Vista questa proprietà si assumerà d'ora in poi che il grafo sia indiretto e che l'adiacenza di un blocco con i vicini sia rappresentata da un solo arco $(blocco,vicino) \equiv (vicino,blocco)$ che sostituirà i due archi, quello entrante e quello uscente, senza alterare la struttura del dominio come si può vedere dalla figura \ref{fig:domain_graph}.\\
% \begin{figure}[H]
% 	\centering
% 	\includegraphics[width=1.0\textwidth]{img/logo_unipr}
% 	\caption{Grafo di un dominio reale senza multi risoluzione}
% 	\label{fig:domain_graph}
% \end{figure}

% \subsubsection{La rappresentazione del grafo}
% La rappresentazione effettiva del grafo su macchina avviene attraverso l'array \textit{neigh} (paragrafo \ref{subsubsec:data_types_swe}) che combina i vantaggi di una rappresentazione a lista e matrice di adiacenza (paragrafo \ref{subsec:graph_rappr}). Questa combinazione è possibile grazie al numero piccolo (otto) di vicini di ogni blocco fissato a priori. Questa proprietà permette all'array \textit{neigh} di poter accedere efficientemente alle informazioni dei vicini di ogni blocco mantenendo comunque un'occupazione di memoria pari a $\Theta(V+E)$.\\
% Ad esempio dato un blocco \textbf{K} è possibile controllare le informazioni di tutti i suoi vicini scorrendo le otto celle da \textit{neigh}[\textbf{K}$.id*8$] a \textit{neigh}[\textbf{K}$.id*8 + 7$]. Inoltre se si vuole controllare lo stato di un vicino in una determinata posizione (Nord, Sud, $\ldots$) basterà seguire lo schema indicato in figura \ref{fig:disp_neigh} al capitolo \ref{sec:first_swe} ed utilizzare l'indice relativo alla posizione voluta: le informazioni sul vicino Nord di \textbf{K}, ad esempio, si possono ricavare, con un singolo accesso, nella cella \textit{neigh}[\textbf{K}.$id*8 + 6$], essendo il vicino Nord associato all'indice 6.

% \subsubsection{La multi risoluzione nel grafo}
% La topologia del grafo che rappresenta il dominio è profondamente influenzata dal concetto di multi risoluzione la cui presenza varia significativamente il rapporto tra la quantità di nodi presenti e l'ampiezza del terreno tra le zone di alto e basso interesse. Un'alta risoluzione crea quindi zone computazionalmente pesanti e nelle quali intercorrono grandi quantità di connessioni mentre, al contrario, zone a basse risoluzioni sono contraddistinte da un esiguo numero di blocchi e da poche comunicazioni.\\
% Questa differenza tra le zone a diverse risoluzioni sarà il punto centrale dell'algoritmo di partizionamento. Questo infatti, cercherà di suddividere il dominio in modo che le zone ad alto costo computazionale siano distribuite equamente tra i vari sotto-domini eseguendo il \textit{taglio} in zone con un basso numero di comunicazioni al fine di ridurre l'\textit{overhead} dovuto allo scambio dei dati.
% \begin{figure}[H]
% 	\centering
% 	\begin{subfigure}{1.0\textwidth}
% 		\centering
% 		\includegraphics[width=1.0\textwidth]{img/logo_unipr}
% 		\caption{Grafo con vertici proporzionali alla zona di terreno rappresentata\newline}
% 		\label{fig:graph_multi}
% 	\end{subfigure}
% 	\begin{subfigure}{1.0\textwidth}
% 		\centering
% 		\includegraphics[width=1.0\textwidth]{img/logo_unipr}
% 		\caption{Grafo con vertici uniformi}
% 		\label{fig:graph_multi_equiv}
% 	\end{subfigure}
% 	\caption{Grafo di un dominio reale con multi risoluzione}
% 	\label{fig:domain_graph_multi}
% \end{figure}

% Il grafo d'ora in avanti sarà rappresentato con vertici rettangolari e senza rappresentare gli archi che connettono i vari nodi. Questa rappresentazione (realizzata tramite l'applicativo Graphviz \cite{graphviz}), mostrata in figura \ref{fig:square_graph} permette una migliore visione d'insieme del dominio
% \begin{figure}[H]
% 	\centering
% 	\includegraphics[width=1.0\linewidth]{img/logo_unipr}
% 	\caption{Modello di rappresentazione utilizzato}
% 	\label{fig:square_graph}
% \end{figure}

% \subsection{Le strutture dati dell'algoritmo}
% Lo sviluppo dell'algoritmo ha richiesto l'introduzione di alcune strutture dati che permettessero di poter elaborare efficientemente il dominio, legando ad ogni nodo del grafo (ovvero ad ogni blocco) informazioni aggiuntive.\\
% Si è quindi sviluppato il tipo di dato \emph{Block} (mostrato nel codice \ref{lst:block_struct}) che permette di salvare al suo interno, oltre alle informazioni fornite dal grafo di partenza (\textit{id\_block, id\_resolution, neighbors}), altri parametri utili all'algoritmo di suddivisione.\\
% \begin{lstlisting}[ 
% label={lst:block_struct},
% caption={Struttura dati \textit{Block}}]
% 	struct Block
% 	{
% 		int id_block;
% 		char resolution;
% 		int id_subdomain;
% 		int key;
% 		bool in_other_subdomain;
% 		struct neigh_t neighbors[4];
% 	};
% \end{lstlisting}
% \textit{Block} cambia leggermente la modalità di memorizzazione dei vicini trasportandone tutte le informazioni nell'array \textit{neighbors} che mantiene comunque il tipo \textit{neigh\_t} (paragrafo \ref{subsubsec:data_types_swe}) per la memorizzazione dei vicini).\\
% Al contrario di quello che ci si aspetterebbe la dimensione di \textit{neighbors} è quattro e non otto. Questo è dovuto al fatto che le comunicazioni del blocco con i vicini diagonali avvengono solo in casi di adiacenza particolari e quindi non influenzano i criteri per la suddivisione. L'associazione tra indice e posizione del blocco viene comunque mantenuta introducendo un nuovo schema mostrato in figura \ref{fig:new_scheme_neigh}.\\
% \begin{figure}[H]
% 	\centering
% 	\includegraphics[width=1.0\textwidth]{img/logo_unipr}
% 	\caption{Il nuovo schema di associazione}
% 	\label{fig:new_scheme_neigh}
% \end{figure}

% Il grafo, durante il partizionamento, sarà rappresentato quindi da un array di \textit{Block}, chiamato \emph{adjacency\_list}, che rispecchia le proprietà di \textit{neigh}.

% \subsection{La suddivisione tramite multi-MST}
% Introdotte le strutture dati principali e spiegata la conformazione del grafo si possono passare a descrivere i primi stadi dell'algoritmo pensato per trovare una buona suddivisione.\\
% Una buona ripartizione del dominio sui vari nodi di calcolo dovrebbe, come detto prima, soddisfare due requisiti:
% \begin{itemize}
% 	\item Bilanciare la quantità di informazioni, e quindi il numero di blocchi, fra tutti i sotto-domini;
% 	\item Ridurre al minimo il numero di comunicazioni tra i nodi di calcolo cercando quindi di eseguire i \textit{tagli} al grafo in zone in cui intercorrono pochi \textit{archi}.
% \end{itemize}

% Considerando la multi risoluzione e la conseguente topologia del grafo i due parametri citati precedentemente possono essere rivisti nei seguenti termini:
% \begin{itemize}
% 	\item Suddividere più equamente possibile le zone a risoluzione più alta;
% 	\item Eseguire i \textit{tagli} nelle zone a risoluzione più bassa.
% \end{itemize}

% L'idea su cui si basa l'algoritmo di partizionamento, che cerca di rispettare questi due requisiti, si fonda sul concetto di \textit{albero di copertura minimo} con archi equi-pesati (paragrafo \ref{subsec:mst_description}).\\
% In particolare si utilizzerà una foresta di MST che esplorano il grafo contemporaneamente.\\


% \subsubsection{Suddivisione di un dominio Ideale}
% Inizialmente l'idea è nata immaginando di effettuare una suddivisione su un \emph{dominio ideale} in cui ci sono tante zone ad alta risoluzione quanti nodi di calcolo e dove queste degradano lentamente verso zone a risoluzione minima. Prendendo questo dominio è facile immaginarne una suddivisione ottimale tramite gli alberi di copertura minimi.\\
% In particolare ogni albero di copertura minima, che rappresenta la partizione da assegnare ad un nodo di calcolo, ingloberebbe tutta una zona a risoluzione massima e i conseguenti contorni a risoluzioni intermedie, fermandosi (come mostrato nella figura \ref{fig:mst_stop}) quando incontra le \emph{foglie} di un altro MST (che avrà a sua volta inglobato nei suoi nodi interni le zone a risoluzione maggiore).\\
% Data la configurazione ideale del dominio si può supporre che l'incontro fra i vari MST avvenga nelle zone di risoluzione minima individuando un'area nella quale è possibile  effettuare un \textit{taglio} che minimizza le comunicazioni.
% \begin{figure}[H]
% 	\centering
% 	\includegraphics[width=0.5\textwidth]{img/logo_unipr}
% 	\caption{Le linee di arresto dei vari MST che compongono la foresta}
% 		\label{fig:mst_stop}
% 	\centering{In ogni vertice del grafo è mostrato il realtivo peso e il peso di ogni arco è considerato unitario}

% \end{figure}
% I problemi relativi a questo primo passo dell'algoritmo sono però evidenti: si richiede un dominio ideale e quindi non compatibile con i casi reali, gli unici che sono realmente interessanti, e inoltre si presuppone di conoscere a priori i blocchi di partenza che permettono all'algoritmo di eseguire una suddivisione ottimale.\\

% \subsubsection{Le radici Casuali}\label{subsubsec:random_roots}
% Il partizionamento tramite \emph{multi-MST} (ovvero l'esecuzione di più MST su uno stesso dominio) rimane però un buon punto di partenza. La suddivisione ottenuta grazie a questo approccio, se scelti buoni punti di partenza, riesce a sfruttare la topologia del grafo per creare partizioni bilanciate e \textit{tagli} con un basso numero di comunicazioni.\\
% Partendo dalle giuste radici, i vari alberi di copertura minimi permetterebbero una divisione omogenea del numero dei blocchi.\\
% Avendo infatti gli archi tutti lo stesso peso, la copertura di un albero si fermerà solo quando le foglie di questo incontreranno le foglie di un altro albero con un peso simile (come mostrato in figura \ref{fig:mst_stop}), indicando che il numero di blocchi inglobati è equilibrato.\\
% Inoltre è possibile supporre che, se le radici fossero poste nei punti centrali delle zone a densità di informazione massimale\footnote{Con risoluzione più alta e con massima distanza dalle zone a risoluzione più bassa}, questo tipo di suddivisione permetterebbe di eseguire i \textit{tagli} nelle zone topologicamente più leggere e quindi dove intercorrono meno comunicazioni.\\
% Idealmente infatti posizionando le radici nei punti \textquotedblleft più interni\textquotedblright (ovvero nei punti al centro delle zone più dense) del dominio, le zone sarebbero mappate sull'albero all'incirca in base alla loro concentrazione di blocchi. In particolare le zone a densità minore rappresenterebbero le foglie dell'albero e quindi le zone in cui viene eseguito il \textit{taglio}.\\
% Come primo stadio dell'algoritmo su casi reali è stata realizzata una suddivisione basata su multi-MST che sceglie casualmente le proprie radici e che viene eseguita ripetutamente per poter scegliere una soluzione accettabile.\\
% Nella (figura \ref{fig:random_root}) è mostrata una suddivisione con radici di partenza casuali.
% \begin{figure}[H]
% 	\centering
% 	\includegraphics[width=0.9\textwidth]{img/logo_unipr}
% 	\caption{Ripartizione del dominio in quattro partendo da radici casuali}
% 		\label{fig:random_root}
% \centering{Le radici sono messe in evidenza con colori diversi}

% \end{figure}
% \paragraph{L'implementazione dell'algoritmo di Prim con multi-radice}~\\
% Nel codice \ref{alg:my_prim_multi} si dà una descrizione in pseudo-codice della versione dell'algoritmo di Prim (algoritmo {\ref{alg:Prim_Mst}), utilizzato per realizzare la suddivisione, che ammette un numero di \emph{radici}, e di conseguenza un numero di MST, parametrico.\\
% \textit{roots} è la lista delle radici e \textit{adjacency\_list} rappresenta l'array di \textit{Block} citato precedentemente.
% \begin{algorithm}[H]
% 	\caption{Prim Multi Source}
% 	\label{alg:my_prim_multi}
% 	\begin{algorithmic}[1]
% 		\Statex
% 		\Function{PRIM\_MULTI-MST}{roots\_list \textit{roots}, array\_Block \textit{adjacency\_list}}
		
% 		\State{\color{blue}{//Inizializzazione di tutti i vertici}}
% 		\For{ogni \textit{u} $\in$ \textit{adjacency\_list}}
% 		\Let{\textit{u.key}}{$\infty$}
% 		\Let{\textit{u.id\_subdomains}}{\textit{NIL}}
% 		\Let{\textit{u.in\_other\_subdomain}}{FALSE} 
% 		\EndFor
		
% 		\State{\color{blue}{//Inizializzazione di tutte le radici}}
% 		\Let{\textit{i}}{0}
% 		\For{ogni \textit{r} $\in$ \textit{roots}}
% 		\Let{\textit{r.key}}{0}
% 		\Let{\textit{r.id\_subdomains}}{\textit{i}}
% 		\Let{\textit{i}}{$i+1$}
% 		\EndFor
% 		\State{\textit{Q}: coda di priorita ordinata in base al campo 
% 			\textit{key}}
% 		\Let{\textit{Q}}{\textit{roots}}
% 		\Let{\textit{count}}{0}
		
% 		\While{count $<$ \textit{adjacency\_list.size}}
% 		\Let{\textit{u}}{EXTRACT-MIN(\textit{Q})}
		
% 		\If{!(\textit{u.in\_other\_subdomain})}
% 		\State{\color{blue}{//Essendo il minimo viene aggiunto definitivamente}}
% 		\Let{\textit{u.in\_other\_subdomain}}{TRUE}
% 		\State{\color{blue}{//Si analizzano i vicini}}
% 		\For{ogni \textit{v} $\in$ \textit{u.neighbors}}
% 		\State{\color{blue}{//In questo caso specifico il peso di ogni arco vale 1}}
% 		\If{$v \in adjacency\_list$ and ($u.key+1) < v.key$}
% 		\Let{\textit{v.key}}{$u.key+1$}
% 		\Let{\textit{v.subdomains}}{\textit{u.subdomains}}
% 		\Let{\textit{count}}{\textit{count} + 1}
% 		\EndIf
% 		\EndFor
% 		\EndIf
% 		\EndWhile
% 		\State \Return{}
% 		\EndFunction
% 	\end{algorithmic}
% \end{algorithm}
% Terminata l'esecuzione, per ricostruire i vari sotto-domini, basterà analizzare i campi \textit{id\_subdomain} di tutti gli elementi (blocchi) contenuti all'interno dell'array \textit{adjacency\_list}, oppure memorizzare, per esempio in un array di liste di interi, i vari \textit{id} relativi ai blocchi che compongono un determinato sotto-dominio ogni volta che un blocco viene inglobato definitivamente da un MST (algoritmo \ref{alg:my_prim_multi}, riga 19).

% \subsection{La selezione delle radici}
% Si è implementata una suddivisione basata sulle capacità degli alberi di copertura minimi di sfruttare la conformazione topologica del grafo con multi risoluzione, che utilizza però come radici vertici casuali, richiedendo così un alto numero di tentativi per ottenere una suddivisione soddisfacente.\\
% Per ovviare al problema quindi, si è cercato un criterio di \emph{selezione delle radici}.\\
% Il criterio che verrà proposto in questa tesi si baserà sul concetto di \emph{massimo locale}.\\
% Prima di introdurre il concetto di massimo locale però è necessario parlare di \emph{distance transform}.

% \subsubsection{Distance Transform e Massimi Locali}\label{subsubsec:local_max_map_transf}
% Il concetto di \textit{distance transform} o \textit{distance map} è solitamente legato alla rappresentazione di immagini e consiste nel trasformare un'immagine binaria\footnote{con solo 2 livelli di colore, tipicamente il bianco e il nero} in un'immagine a scala di grigi \cite{distance_transform}. Questa trasformazione mantiene invariata la struttura dell'immagine ma \textquotedblleft sfuma\textquotedblright~il livello di colore partendo dai punti più interni ed arrivando, con i toni più chiari, ai bordi dell'immagine.
% \begin{figure}[H]
% 	\centering
% 	\begin{subfigure}{0.5\textwidth}
% 	\centering
% 	\includegraphics[width=0.5\textwidth]{img/logo_unipr}
% 	\caption{L'immagine prima della trasformazione}
% 	\end{subfigure}%
% 	\begin{subfigure}{0.5\textwidth}
% 	\centering
% 	\includegraphics[width=0.5\textwidth]{img/logo_unipr}
% 	\caption{La stessa immagine dopo la trasformazione}
%     \end{subfigure}
% 	\caption{Esempio di \textit{distance transform} su immagine binaria}
% 	\label{fig:map_transform_example}
% \end{figure}

% Lo stesso concetto può essere applicato al contrario, ovvero: dati i punti più esterni, si può \textquotedblleft sfumare\textquotedblright~verso quelli più interni cioè verso tutti quei punti con una massima distanza dai bordi rispetto ai loro vicini.\\
% Una suddivisione di questo genere, applicata al dominio che si intende partizionare, permette di creare una distribuzione discreta e lineare di pesi.\\
% In particolare si può utilizzare una specifica trasformazione \textit{distance map} che, partendo dai punti centrali delle zone a risoluzione minima locale (che comunicano cioè solo con zone a risoluzione maggiore), distribuisce i valori aumentandoli man mano che ci si allontana.\\
% Dopo questa trasformazione si avrà quindi una \textquotedblleft curva\textquotedblright~della distribuzione delle risoluzioni ovvero, si passerà da una distribuzione di valori quantizzati\footnote{I valori rappresentano i vari livelli di risoluzione e sono: 1, 2, 3 e 4} molto distanti, ad una distribuzione più fitta. Questa distribuzione considera le distanze dalle risoluzioni vicine in modo da poter guidare una ricerca per gradiente dei massimi locali (figura \ref{fig:local_max_graph}).
% I vertici del grafo che avranno associato un valore maggiore rispetto ai propri vicini, saranno identificati proprio come \emph{massimi locali} e rappresenteranno i punti più interni delle zone che hanno localmente la massima risoluzione.\\
% \begin{figure}[H]
% 	\centering
% 	\includegraphics[width=1.0\textwidth]{img/logo_unipr}
% 	\caption{Grafo del dominio con i massimi locali evidenziati (in nero)}
% 	\label{fig:local_max_graph}
% \end{figure}

% Si illustra di seguito lo pseudo codice relativo alla trasformazione utilizzata sul dominio (che sfrutta l'algoritmo \ref{alg:my_prim_multi}) per l'individuazione dei massimi locali:
% \begin{algorithm}[H]
% 	\caption{Map Tranform}
% 	\label{alg:my_map_transfor}
% 	\begin{algorithmic}[1]
% 		\Statex
% 		\Function{map\_tranform}{array\_Block \textit{adjacency\_list}}
% 		\State{\color{blue}{//Inizializzazione di tutti i vertici}}
% 		\For{ogni \textit{u} $\in$ \textit{adjacency\_list}}
% 		\Let{\textit{u.key}}{$\infty$}
% 		\Let{\textit{u.in\_other\_subdomain}}{FALSE} 
% 		\EndFor
		
% 		\State{\color{blue}{//Si analizza il grafo un livello di risoluzione alla volta}}		\Let{\textit{res}}{\textit{min\_res\_value}}
% 		\While{\textit{res} $\leq$ \textit{max\_res\_value}}
% 		\State{\color{blue}{Si definisce \textit{adjacency\_list}(\textit{res\_value}) l'insieme}}
% 		\State{\color{blue}{dei nodi a risoluzione \textit{res\_value}}}	
% 		\Let{\textit{Q}}{\textit{adjacency\_list}(\textit{res\_value})}
% 		\Let{is\_L}{FALSE}
% 		\While{{Q} $\neq 0$}
% 		\If{EXISTS\_LOW(\textit{Q})\color{orange}{$^1$}\color{black} and \textit{res} $\neq$ \textit{min\_res\_value}}
% 		\Let{\textit{start\_node}}{EXTRACT\_BOUNDARY\_L(\textit{Q})\color{orange}{$^2$}\color{black}}
% 		\Let{\textit{is\_L}}{TRUE}
% 		\Else
% 		\Let{\textit{start\_node}}{EXTRACT\_BOUNDARY\_H(\textit{Q})\color{orange}{$^3$}\color{black}}
% 		\EndIf
% 		\State{MST\_WEIGHT\_DIST(\textit{start\_node, is\_L, adjacency\_list, Q})\color{orange}{$^4$}}
% 		\EndWhile
% 		\EndWhile
% 		\State \Return{}
% 		\EndFunction
% 	\end{algorithmic}
% \end{algorithm}
% \begin{itemize}
% \color{orange}\item[1]\color{black}EXISTS\_LOW (riga 10) verifica se è presente almeno un blocco che ha almeno un vicino a risoluzione minore;
% \color{orange}\item[2]\color{black}EXTRACT\_BOUNDARY\_L (riga 11) estrae come blocco di partenza un blocco che comunica con vicini a risoluzione minore;
% \color{orange}\item[3]\color{black}EXTRACT\_BOUNDARY\_H (riga 14) estrae un blocco che sia sul bordo del dominio o che abbia almeno un vicino a risoluzione maggiore;
% \color{orange}\item[4]\color{black}MST\_WEIGHT\_DIST (riga 15) spiegato nell'algoritmo \ref{alg:dist_weight_mst}.
% \end{itemize}

% \begin{algorithm}[H]
% 	\caption{MST Weight Redistribution}
% 	\label{alg:dist_weight_mst}
% 	\begin{algorithmic}[1]
% 		\Statex
% 		\Function{MST\_WEIGHT\_DIST}{Block \textit{start\_node}, bool \textit{is\_L}, array\_Block \textit{adjacency\_list}, list\_same\_res\ \textit{Q}}
% 	    \Let{T}{EXTRACT\_COMMUNICANT(\textit{Q}, \textit{res\_start})\color{mygreen}{$^1$}\color{black}}
		
% 		\If{\textit{is\_L}}
	   
% 		\Let{\textit{list\_start\_node}}{COMPLETE\_BOUNDARY\_L(\textit{start\_node}, \textit{T})\color{mygreen}{$^2$}\color{black}}
		
% 		\Let{\textit{new\_start}}{PRIM\_MULTI-MST\_central(\textit{list\_start\_node}, \textit{T})\color{mygreen}{$^3$}\color{black}}
% 		\Let{\textit{T}[\textit{new\_start}]\textit{.key}}{0}
% 		\Else
		
	
% 		\Let{\textit{list\_start\_node}}{COMPLETE\_BOUNDARY\_H(\textit{start\_node}, \textit{T})\color{mygreen}{$^4$}\color{black}}
% 		\State{BOUNDARY\_WEIGHT(\textit{list\_start\_node}, \textit{adjacency\_list})\color{mygreen}{$^5$}\color{black}}
% 		\EndIf
% 		\State{PRIM\_MULTI-MST\_final(\textit{list\_start\_node}, \textit{T}, \textit{Q}, \textit{adjacency\_list})\color{mygreen}{$^6$}\color{black}}
% 		\State \Return{}
% 		\EndFunction
% 	\end{algorithmic}
% \end{algorithm}
% \begin{itemize}
% \color{mygreen}\item[1]\color{black} EXTRACT\_COMMUNICANT (riga 2) estrae tutti i nodi di \textit{Q} che hanno comunicazione (diretta o indiretta) con \textit{start\_node} passando solo attraverso blocchi della stessa risoluzione;
% \color{mygreen}\item[2]\color{black} COMPLETE\_BOUNDARY\_L (riga 4) estrae tutti i nodi nella zona di \textit{start\_node} (\textit{T}) che hanno stessa risoluzione e che possiedono almeno un vicino a risoluzione minore;
% \color{mygreen}\item[3]\color{black} Si considera che la funzione PRIM\_MULTI-MST\_central (riga 5) esegua l'algoritmo \ref{alg:my_prim_multi} ma ritorni il nodo con valore \textit{key} maggiore senza modificare globalmente \textit{T};
% \color{mygreen}\item[4]\color{black} COMPLETE\_BOUNDARY\_H (riga 8) estrae tutti i nodi nella zona di \textit{start\_node} (\textit{T}) che hanno stessa risoluzione e che possiedono almeno un vicino a risoluzione maggiore;
% \color{mygreen}\item[5]\color{black}BOUNDARY\_WEIGHT (riga 9) ridefinisce i pesi delle radici (\textit{list\_start\_node}) assegnando ad ogni radice il peso del proprio vicino (anche a risoluzione minore) \textquotedblleft meno pesante\textquotedblright aumentato di 1 (tutti gli archi hanno infatti peso unitario);
% \color{mygreen}\item[6]\color{black} PRIM\_MULTI-MST\_final (riga 10) esegue l'algoritmo \ref{alg:my_prim_multi} su \textit{T} senza settare a 0 il valore \textit{key} delle radici. Inoltre estrae da \textit{Q} tutti i blocchi in \textit{T} e riporta le modifiche sui valori dei blocchi in \textit{adjacency\_list}.
% \end{itemize}


% \subsubsection{I massimi locali randomici come partenza}
% Grazie all'algoritmo \ref{alg:my_map_transfor} è quindi possibile determinare tutti i massimi locali del grafo; essendo questi i punti più interni del grafo, come detto nel paragrafo \ref{subsubsec:random_roots}, costituiscono delle buone radici per i vari MST.\\
% Siccome il numero di partizioni desiderate è quasi sempre di molto inferiore al numero di massimi locali trovati, nasce il problema di quali scegliere come punti di partenza.\\
% Un primo approccio è stato quello di scegliere le radici dell'algoritmo di partizione (algoritmo \ref{alg:my_prim_multi}) casualmente tra tutti i massimi locali (figura \ref{fig:randomo_local_max}).\\
% Come si nota in figura \ref{fig:randomo_local_max} però questo approccio non garantisce affatto una suddivisione ideale.
% \begin{figure}[H]
% 	\centering
% 	\includegraphics[width=1.0\textwidth]{img/logo_unipr}
% 	\caption{Ripartizione del dominio in quattro partendo da radici scelte casualmente tra i massimi locali}
% 	\label{fig:randomo_local_max}
% \end{figure}


% \subsection{Multi-MST sui massimi locali}
% Calcolati tutti i massimi locali bisogna quindi cercare di sfruttarli per produrre una buona suddivisione.\\
% I massimi locali sono, come definito nel paragrafo \ref{subsubsec:local_max_map_transf}, i punti centrali delle zone a densità di informazione massimale (risoluzione più alta e lontana da zone a risoluzione più bassa) e quindi delle zone più dense del grafo. Vista la posizione dei massimi locali è facile immaginare che nelle loro vicinanze vi sia un'alta concentrazione di blocchi e, di conseguenza un alto numero di comunicazioni.\\
% Se si prende, per esempio, il dominio in figura \ref{fig:local_max_graph} nel quale sono stati calcolati i massimi locali risulta infatti evidente che le zone dove intercorrono meno comunicazioni sono quelle più distanti dai vari massimi locali.\\
% Questo è dovuto proprio al fatto che i massimi locali sono al centro delle zone più pesanti dal grafo e quindi allontanandosi da questi si arriva su zone con una bassa densità di blocchi (zone a bassa risoluzione).\\

% \subsubsection{Zone di appartenenza e Tagli ideali}\label{subsubsec:zone_app}
% Vista questa proprietà si può quindi supporre che sia conveniente effettuare i \textit{tagli} tra i blocchi topologicamente più distanti dai vari massimi locali.\\
% Per individuare queste zone si è sfruttata ancora una volta l'idea di albero di copertura minimo.\\
% In particolare si è pensato di sviluppare contemporaneamente, sullo stesso dominio, un MST per ogni massimo locale (con radice nel massimo locale stesso) utilizzano l'algoritmo \ref{alg:my_prim_multi}.\\
% Quando l'algoritmo ha terminato la sue esecuzione si potranno distinguere nel grafo le varie \textquotedblleft \emph{zone di appartenenza}\textquotedblright, una per ogni massimo locale (figura \ref{fig:mst_local_max}).
% \begin{figure}[H]
% 	\centering
% 	\includegraphics[width=1.0\textwidth]{Immagini/zone_di_appartenenza.png}
% 	\caption{Le zone di appartenenza descritte dalla foresta di MST}
% 	\label{fig:mst_local_max}
% \end{figure}

% Una volta definite le \textquotedblleft zone di appartenenza\textquotedblright~risulta quindi chiaro che i bordi esterni dei vari MST rappresentano le zone più distanti dai massimi locali.\\
% In particolare, seguendo l'idea che nelle nelle zone più lontane dai massimi locali ci sia minima comunicazione, i vari \textit{tagli} dovrebbero essere effettuati lungo \emph{le linee di contatto} (mostrate in figura \ref{fig:mst_stop}) tra i vari MST.\\
% Tagliando lungo le suddette linee infatti è possibile evitare di \textit{tagliare} dove vi è una grande densità di blocchi.\\
% Si può dire quindi che le \textit{linee di contatto} forniscano le zone ideali su cui effettuare i vari \textit{tagli} necessari a partizionare il grafo.

% \subsubsection{Unione delle zone di appartenenza}\label{subsubsec:varius_merge}
% Definite le linee lungo le quali effettuare i possibili \textit{tagli} è il momento di suddividere il grafo vero e proprio nel numero di partizioni desiderato.\\
% Per mantenere intatte le zone di taglio sicure si è pensato di eseguire la suddivisione partendo proprio dalle \textit{zone di appartenenza}.\\
% Unendo tra loro le singole zone è possibile infatti creare sotto-domini maggiori che conservino le \textit{linee di taglio} individuate precedentemente e che producano quindi un buon partizionamento del grafo.\\
% Perciò si è usato un processo di unione iterativo che ad ogni passo seleziona i due sotto-domini migliori da unire fino a raggiunger il numero di partizioni volute.\\
% Siccome vi sono diversi criteri di unione validi, questo algoritmo di unione è stato sviluppato in più versioni ognuna delle quali può variare sia nella scelta iniziale del sotto-dominio che si vuole unire, sia nella scelta della partizione da unire con il suddetto sotto-dominio.\\
% In particolare si avranno i seguenti criteri:
% \begin{itemize}
% 	\item[1] Si sceglie come sotto-dominio di partenza quello con il rapporto\\
% 	$numero\_blocchi\mathbin{/}numero\_comunicazioni$ peggiore (cioè minore) e come partizione da unire quella che riduce maggiormente le comunicazioni del sotto-dominio di partenza (figura \ref{fig:nb_comm});
% 	\item[2] Si sceglie come sotto-dominio di partenza quello con il rapporto\\
% 	$numero\_blocchi\mathbin{/}numero\_comunicazioni$ peggiore (cioè minore) e come partizione da unire quella che, unita al sotto-dominio di partenza, produce il rapporto $numero\_blocchi\mathbin{/}numero\_comunicazioni$ migliore (figura \ref{fig:nb_nb});
% 	\item[3] Si sceglie come sotto-dominio di partenza quello con il più alto numero di comunicazioni e come partizione da unire quella che,che riduce maggiormente le comunicazioni del sotto-dominio di partenza (figura \ref{fig:comm_comm}).
% \end{itemize}
% Tutti e tre i vari criteri sono suddivisi in altre due versioni: una che ammette la possibilità di eseguire un'unione anche tra vicini non adiacenti per permettere un bilanciamento migliore, e un'altra che non lo consente (a meno di un grande sbilanciamento nei sotto-domini finali).\\
% \begin{figure}[H]
% 	\centering
% 	\begin{subfigure}{1.0\textwidth}
% 	\centering
% 		\begin{subfigure}{0.5\textwidth}
% 			\centering
% 			\includegraphics[width=0.9\textwidth]{img/logo_unipr}
% 			Amessi solo vicini adiacenti
% 		\end{subfigure}%
% 		\begin{subfigure}{0.5\textwidth}
% 			\centering
% 			\includegraphics[width=0.9\textwidth]{img/logo_unipr}
% 			Vicini non adiacenti ammessi
% 		\end{subfigure}
% 	\caption{Ottimizzazione rapporto $numero\_blocchi\mathbin{/}numero\_comunicazioni$ e comunicazioni\newline}
% 	\label{fig:nb_comm}
% 	\end{subfigure}
% 	\begin{subfigure}{1.0\textwidth}
% 	\centering
% 		\begin{subfigure}{0.5\textwidth}
% 			\centering
% 			\includegraphics[width=0.9\textwidth]{img/logo_unipr}
% 			Amessi solo vicini adiacenti
% 		\end{subfigure}%
% 		\begin{subfigure}{0.5\textwidth}
% 			\centering
% 			\includegraphics[width=0.9\textwidth]{img/logo_unipr}
% 			Vicini non adiacenti ammessi
% 		\end{subfigure}
% 	\caption{Ottimizzazione rapporto $numero\_blocchi\mathbin{/}numero\_comunicazioni$\newline}
% 	\label{fig:nb_nb}
% 	\end{subfigure}
% 	\begin{subfigure}{1.0\textwidth}
% 	\centering
% 	\begin{subfigure}{0.5\textwidth}
% 		\centering
% 		\includegraphics[width=0.9\textwidth]{img/logo_unipr}
% 		Amessi solo vicini adiacenti
% 	\end{subfigure}%
% 	\begin{subfigure}{0.5\textwidth}
% 		\centering
% 		\includegraphics[width=0.9\textwidth]{img/logo_unipr}
% 		Vicini non adiacenti ammessi
% 	\end{subfigure}
% 	\caption{Ottimizzazione comunicazioni\newline}
% 	\label{fig:comm_comm}
% \end{subfigure}
% 	\caption{I vari criteri di unione a confronto}
% \end{figure}

% I vari criteri di unione creano diverse suddivisioni all'interno dello stesso grafo ma producono comunque delle ripartizioni accettabili, sia dal punto di vista del bilanciamento che dal punto di vista dell comunicazioni, senza richiedere diversi tentativi come invece accade nelle implementazioni basate su radici casuali.\\
% Si è quindi riusciti ad ottenere una suddivisione soddisfacente del dominio, che permette di ripartire il carico adeguatamente su una struttura MPI.


% \subsection{I confronti}
% Nell'immagine \ref{fig:comm_confronto} si mostreranno due grafici che valutano in base a due parametri (il bilanciamento e il numero delle comunicazioni totali) le varie fasi dell'algoritmo e i vari criteri di unione delle \textit{zone di appartenenza} (paragrafo \ref{subsubsec:varius_merge}).\\
% In particolare si sono messi a confronto i vari criteri di unione, tra loro e con altri algoritmi di suddivisione basati sulla posizione dei blocchi o su scelte di radici casuali.
% %\begin{itemize}
% %	\item [1] Prim multi-MST con unione che ottimizza il numero delle comunicazioni;
% %	\item [2] Prim multi-MST con unione che ottimizza il numero delle comunicazioni consentendo anche unione tra sotto-domini non adiacenti;
% %	\item [3] Prim multi-MST con unione che ottimizza il rapporto $numero\_blocchi\mathbin{/}numero\_comunicazioni$ e il numero di comunicazioni;
% %	\item [4] Prim multi-MST con unione che ottimizza il rapporto $numero\_blocchi\mathbin{/}numero\_comunicazioni$ e il numero di comunicazioni consentendo anche unione tra sotto-domini non adiacenti;
% %	\item [5] Prim multi-MST con unione che ottimizza il rapporto $numero\_blocchi\mathbin{/}numero\_comunicazioni$;
% %	\item [6] Prim multi-MST con unione che ottimizza il rapporto $numero\_blocchi\mathbin{/}numero\_comunicazioni$ consentendo anche unione tra sotto-domini non adiacenti;
% %	\item [7] Suddivisione del dominio suddividendo in parti uguali l'asse \textit{x};
% %	\item [8] Suddivisione del dominio suddividendo in parti uguali l'asse \textit{y};
% %	\item [9] Suddivisione tramite multi-MST che utilizza come punti di partenza posizioni casuali (risultato migliore su 50 tentativi);
% %	\item [10] Suddivisione tramite multi-MST che utilizza come punti di partenza \textit{massimi locali} casuali (risultato migliore su 50 tentativi);	
% %\end{itemize}
% \begin{figure}[H]
% 	\centering
% 	\begin{subfigure}{1.0\textwidth}
% 		\centering
% 		\includegraphics[width=0.9\textwidth]{img/logo_unipr}
% 		\caption{Il confronto degli algoritmi con quattro partizioni\newline}
% 		\label{fig:confronto_4}
% 	\end{subfigure}
% 	\begin{subfigure}{1.0\textwidth}
% 		\centering
% 		\includegraphics[width=0.9\textwidth]{img/logo_unipr}
% 		\caption{Il confronto degli algoritmi con otto partizioni}
% 		\label{fig:confronto_8}
% 	\end{subfigure}

% 	\caption{Le varie suddivisioni a confronto}
% 		\label{fig:comm_confronto}
% \end{figure}


% Nei grafici in figura \ref{fig:comm_confronto} non emerge, come ci si aspetterebbe, un miglioramento sul numero delle comunicazioni tra le suddivisioni basate sulle zone di appartenenza rispetto a quelle in cui si sono scelte radici casuali. Si nota anzi, che la suddivisione basata sull'offset ottimizza il numero di comunicazioni: questo è probabilmente dovuto alla conformazione del dominio e non è attribuibile ad una vera e proprio caratteristica della suddivisione.\\
% Per quanto riguarda il bilanciamento invece, le suddivisioni basate sulle zone di appartenenza, ottengono dei buoni risultati rispetto alle altre. In particolare, si ottengono dei buoni risultati quando si ammettono anche unioni tra sotto-domini non adiacenti.

% \subsubsection{Il risultato finale}
% Il risultato dei confronti tra i vari algoritmi di ripartizione (figura \ref{fig:comm_confronto}), non fa emergere, come ci si aspetterebbe, un netto miglioramento sul numero di comunicazioni dovuto alla tecnica di suddivisione basata sulle \textit{zone di appartenenza}.\\
% Questo è dovuto al fatto che il criterio di ripartizione basato sulle \textit{zone di appartenenza}, suddivide il dominio basando la sua esecuzione su zone che non cambiano mai la loro forma di base; infatti ogni sotto-dominio finale è semplicemente l'unione dei diversi MST iniziali.\\
% Lo scarso miglioramento, rispetto ad algoritmi che si basano su scelte casuali, è quindi probabilmente dovuto alla struttura non ideale del \textit{taglio}, dettata dalla rigidità dei sotto-domini di base. Si può concludere che l'algoritmo basato sulle \textit{zone di appartenenza}, riesca ad individuare dove sia più corretto eseguire i \textit{tagli}, ma non li esegua nel modo più adeguato.\\

% Al contrario, l'ultima suddivisione proposta in questa tesi, garantisce sempre un buon bilanciamento dei blocchi tra le varie ripartizioni.\\
% Questa proprietà è molto importante in quanto, anche se non se ne ha la certezza, la distribuzione non equilibrata del carico porterebbe ad un overhead molto alto. Infatti, anche se non è ancora stata sviluppata una versione che permetta di testare realmente le suddivisioni, sembra che lo squilibrio del carico sui vari nodi causi un maggiore rallentamento rispetto a quello dovuto ad un alto numero di comunicazioni.\\

% L'algoritmo studiato, nella sua ultima versione (paragrafo \ref{subsubsec:zone_app}), ha l'interessante proprietà di individuare un partizionamento con un numero di comunicazioni accettabile e di garantire un bilanciamento quasi perfetto\footnote{Il bilanciamento è migliore se si permette di unire tra loro sotto-domini non adiacenti}.\\
% L'idea di ripartizione proposta in questa tesi potrebbe essere quindi, con qualche piccolo aggiustamento, una buona base per produrre una suddivisione ottima su tutti i vari domini che l'applicativo \textit{sweGPU} deve elaborare, indipendentemente dalla loro confermazione specifica.\\
% In particolare l'idea di partizionamento basata sull'esecuzione di una foresta di MST (con radici adeguate) su uno stesso dominio con multi risoluzione, sembra garantire una suddivisione migliore di quelle che verrebbero individuate con gli algoritmi di partizionamento classici\cite{krapis_lab}.
% In particolare l'algoritmo che ottimizza il rapporto $numero\_blocchi\mathbin{/}numero\_comunicazioni$ e le comunicazioni, e che non ammette i sotto-domini non adiacenti, è quello che produce i risultati migliori.
\chapter{Testing}
% % !TeX spellcheck = it_IT

% \section{Il modello MPI}
% Una volta definito come partizionare il dominio (capitolo \ref{sec:fourth_sudd_dom}) è necessario predisporre l'ambiente di calcolo per eseguire l'applicativo parallelamente sui vari nodi.\\
% Di seguito si introduce un modello MPI che assegna le diverse partizioni sui vari nodi di calcolo e che gestisce le varie comunicazioni tra i nodi stessi.

% \subsection{L'architettura e la distribuzione dei dati}
% Il tipo di architettura hardware pensata per supportare la versione parallelizzata di questo applicativo è un'architettura \textit{MIMD} (figura \ref{fig:mimd}) a memoria distribuita.\\
% In particolare si parla di memoria distribuita in quanto le informazioni della partizione sono note solo al nodo che la contiene e possono essere lette dagli altri solo attraverso uno scambio di messaggi.\\
% Prima che il programma distribuito venga eseguito è quindi necessario assegnare ad ogni nodo le informazioni che gli appartengono e quelle che gli sono utili per eseguire correttamente le comunicazioni con gli altri nodi di calcolo.
% Ogni nodo ha bisogno di conoscere solamente:
% \begin{itemize}
% 	\item Quale partizione è assegnata ad ogni nodo di calcolo;
% 	\item I dati della mappatura del terreno di ogni blocco contenuto nella propria partizione;
% 	\item La lista, identificata dall'array \textit{neigh} (paragrafo \ref{subsubsec:data_types_swe}), di tutti i blocchi con i relativi vicini;
% 	\item La partizione (e di conseguenza il nodo) di appartenenza di ogni blocco del dominio.
% \end{itemize}

% \subsection{La fase di pre-computazione}
% Nel modello MPI avverrà quindi una fase di pre-computazione che si occuperà di distribuire sui vari nodi tutte le informazioni necessarie.\\
% Come tipicamente avviene nelle strutture di calcolo distribuito, la fase di pre-computazione viene eseguita soltanto su un nodo, d'ora in poi indicato con il termine \emph{root}, che si occupa di calcolare e distribuire le varie informazioni.\\
% In particolare, durante la fase di pre-computazione per la parallelizzazione in un ambiente MPI del programma \textit{sweGPU}, il nodo \textit{root} sarà l'unico a conoscere le informazioni sul domino.
% Nello pseudo-codice \ref{alg:pre_compt_mpi} si descrive la fase di pre-computazione (mostrata in figura \ref{fig:pre_proc_mpi}) eseguita dal nodo \textit{root}, considerando \textit{D} come l'insieme delle informazioni riguardanti il dominio e \textit{par\_id\_to\_node} come una funzione che associa l'id di una partizione al rispettivo nodo.
% \begin{algorithm}[H]
% 	\caption{Fase di pre-computazione}
% 	\label{alg:pre_compt_mpi}
% 	\begin{algorithmic}[1]
% 		\Statex
% 		\Function{Data\_Distribuition}{program\_data \textit{D}, func \textit{par\_id\_to\_node}}
% 		\State{MPI::Init()}
% 		\State{\color{blue}{//Invio delle partizioni ai nodi relativi}}
% 		\If{\textit{rank} $\equiv$ \textit{root.rank}}
% 		\Let{\textit{lists\_partitions}}{MAKE\_PARTITION(\textit{D.domain})}
% 		\For{ogni \textit{partition} $\in$ \textit{lists\_partitions}}
% 		\Let{\textit{node\_to\_send}}{\textit{par\_id\_to\_node(partition.id)}}
% 		\If{\textit{node\_to\_send}$\neq$\textit{root.rank}}
% 		\State{MPI\_NON\_BLOCKING\_SEND(\textit{partition}, node\_to\_send)}
% 		\State{\textit{add\_to\_list\_block-partition}(partition)}
% 		\Else
% 		\State{SAVE(\textit{partition})}
% 		\EndIf
% 		\EndFor
% 		\Else
% 		\State{MPI\_RECEIVE(\textit{root}, \textit{partition})}
% 		\State{SAVE({\textit{partition}})}
% 		\EndIf
% 		\State{\color{blue}{/*Invio in \textit{broadcast} (paragrafo \ref{subsubsec:coll_comm}) dell'array \textit{neigh} e dell'array}}
% 		\State{\color{blue}{\textit{block-partition} che associa ad ogni blocco la propria partizione*/}}
% 		\If{\textit{rank} $\equiv$ \textit{root.rank}}
% 		\State{MPI\_BROADCAST\_SEND(\textit{D.neigh})}
% 		\State{SAVE({\textit{D.neigh}})}
% 		\State{MPI\_BROADCAST\_SEND(\textit{block-partition})}
% 		\State{SAVE({\textit{block-partition}})}
% 		\State{DELETE(\textit{D})}
% 		\Else
% 		\State{MPI\_BROADCAST\_RECEIVE(\textit{neigh})}
% 		\State{SAVE({\textit{neigh}})}
% 		\State{MPI\_BROADCAST\_RECEIVE(\textit{block-partition})}
% 		\State{SAVE({\textit{block-partition}})}
% 		\EndIf
% 		\State{MPI::Finalize()}
% 		\State \Return{}
% 		\EndFunction
% 	\end{algorithmic}
% \end{algorithm}

% \subsection{Lo scambio dei Dati di esecuzione}
% Finita la fase di pre-computazione, l'applicativo per la simulazione numerica dei flussi può iniziare ad eseguire i vari calcoli parallelamente sui vari nodi.\\
% Come spiegato nel paragrafo \ref{subsec:sweworking}, durante ogni passo di esecuzione del programma ogni cella ha bisogno delle informazioni dei propri vicini nel dominio reale. In particolare può capitare che una cella di bordo richieda informazioni che si trovano su un altro blocco (figura \ref{fig:cella_bordo_comunicazione}) e siccome il dominio è stato distribuito su vari nodi può capitare che il blocco in questione non si trovi sullo stesso nodo di quello che contiene la cella.
% Quando si verifica questa situazione il nodo deve richiedere, tramite protocolli MPI, le informazioni necessarie ai nodi che contengono i vicini.\\
% Ad ogni passo di calcolo i nodi devono scambiare informazioni riguardanti i bordi tra blocchi adiacenti e su nodi diversi. Per sfruttare la banda di trasmissione, conviene accorpare i dati da scambiare tra coppie di nodi in un unico blocco di informazione.\\
% In particolare si accodano tutte le informazioni dei vari bordi da trasmettere. Ciascun nodo ha informazioni sufficienti per sapere quali bordi dovranno essere spediti e ricevuti. In questo modo alla ricezione del blocco di dati, sarà possibile gestire i singoli bordi.\\

% Nello pseudo-codice \ref{alg:mpi_exchange} viene descritto il modello MPI di comunicazione tra i nodi durante lo svolgimento del programma \textit{sweGPU}, in cui la funzione \textit{par\_id\_to\_node}, associa l'id di una partizione al rispettivo nodo e il parametro \textit{PRE-DATA} contiene le informazioni derivanti dalla pre-computazione.\\
% Inoltre in questo modello vengono eseguite tutte le \textit{send non bloccanti} prima delle \textit{receive} in modo da poter inviare tutti i dati senza attendere la sincronizzazione tra i nodi.\\
% Uno schema del modello è mostrato in figura \ref{fig:my_mpi_model}.
% \begin{algorithm}[H]
% 	\caption{Modello dello scambio di informazioni}
% 	\label{alg:mpi_exchange}
% 	\begin{algorithmic}[1]
% 		\Statex
% 		\Function{sweGPU\_MPI}{func \textit{par\_id\_to\_node}, pre\_data \textit{PRE-DATA}}
% 		\State{MPI::Init()}
% 		\For{ogni passo di computazione del programma \textit{sweGPU}}
% 		\State{\color{blue}{//Calcolo dei dati da inviare ad ogni nodo}}
% 		\Let{\textit{partition}}{\textit{PRE-DATA.partition}}
% 		\For{ogni \textit{block} $\in$ \textit{partition}}
% 		\For{ogni \textit{single\_neigh} $\in$ \textit{PRE-DATA.neigh(block)}}
% 		\If{\textit{single\_neigh.id\_partition} $\neq$ \textit{partition.id}}
% 		\Let{\textit{info\_to\_send}}{\textit{PRE-DATA.block.info}}
% 		\Let{\textit{node\_to\_send}}{\textit{par\_id\_to\_node}(\textit{partition.id})}
% 		\State{push(\textit{list\_data\_to\_send}(\textit{node\_to\_send}), \textit{info\_to\_send})}
% 		\EndIf
% 		\EndFor
% 		\EndFor
% 		\For{ogni altro nodo \textit{node} che ha almeno una comunicazione}
% 		\State{MPI\_NON\_BLOCKING\_SEND(\textit{list\_data\_to\_send}(\textit{node}))}
% 		\EndFor
% 		\For{ogni altro nodo \textit{node} che ha almeno una comunicazione}
% 		\State{MPI\_RECEIVE(node ,\textit{list\_data\_to\_receive})}
% 		\State{SAVE(\textit{list\_data\_to\_receive})}
% 		\EndFor
% 		\State{DO\_THE\_REAL\_JOB\_ON\_MY\_PARTITION}
% 		\State{\color{blue}{//Le informazioni contenute nei vari blocchi cambiano}}
% 		\EndFor
% 		\State{MPI::Finalize()}
% 		\State \Return{}
% 		\EndFunction
% 	\end{algorithmic}
% \end{algorithm}

% Essendo l'\textit{overhead} per la ricezione dei bordi (riga 15 algoritmo \ref{alg:mpi_exchange}) molto penalizzante, si può pensare di ridurlo introducendo la possibilità di eseguire il calcolo durante il tempo di ricezione per i blocchi che non dipendono dai dati attesi .
% \begin{figure}[H]
% \begin{subfigure}{0.5\textwidth}
% 	\centering
%     \includegraphics[width=1.0\textwidth]{img/logo_unipr}
%     \caption{Schema delle comunicazioni nella fase di pre-computazione}
% 	\label{fig:pre_proc_mpi}
% \end{subfigure}%
% \begin{subfigure}{0.5\textwidth}
% 	\centering
% 	\includegraphics[width=1.0\textwidth]{img/logo_unipr}
% 	\caption{Schema delle comunicazioni del modello di scambio informazioni}
% 	\label{fig:my_mpi_model}
% \end{subfigure}
% \caption{La rappresentazione del modello MPI}
% \end{figure}
%\chapter{Conclusioni}
\section{Sviluppi futuri}
% % !TeX spellcheck = it_IT

% \section{Lavori Futuri}

% \subsection*{Variazione sui criteri di suddivisione}
% L'algoritmo di ripartizione studiato in questa tesi si prefigge di giungere ad una suddivisione accettabile del dominio per poterlo distribuire sui vari nodi di calcolo. In particolare non si è ricercato un algoritmo ottimale ma uno che permetta, con un overhead relativamente basso, di eseguire una suddivisione accettabile in base al bilanciamento delle partizioni e al numero di comunicazioni tra i vari nodi.\\
% Non essendo infatti ancora implementata una versione distribuita di \textit{sweGPU} non si conoscono i fattori penalizzanti della distribuzione, in particolare non si conosce l'overhead introdotto dalle comunicazioni e quello introdotto da un eventuale sbilanciamento.\\
% Quando verrà implementata una versione per architetture multi-nodo si potranno valutare gli overhead dei diversi fattori e quindi ridefinire l'algoritmo in base ai vari parametri che maggiormente influenzano la suddivisione stessa. Se, per esempio, si noterà che l'overhead introdotto dalle comunicazioni è insignificante rispetto a quello introdotto da uno sbilanciamento si potrebbe pensare di forzare l'algoritmo a suddividere il numero di celle perfettamente tra i vari nodi anche a se questo dovesse introdurre un alto numero di comunicazioni.
% \subsection*{Migliorie all'algoritmo di partizionamento}
% Vi sono poi alcune migliorie da apportare all'algoritmo analizzato nel capitolo \ref{sec:fourth_sudd_dom}. In particolare:
% \begin{itemize}
% 	\item Poiché l'ultimo stadio dell'algoritmo, per formare le varie suddivisioni, lavora sull'unione di sotto-domini più piccoli, i vari \textit{tagli} che dividono le varie partizioni possono avere forme abbastanza irregolari. Questa irregolarità può portare ad un numero di comunicazioni leggermente maggiore di quanto ci si aspetti.\\
% 	Bisognerebbe quindi implementare una piccola procedura che al termine dell'esecuzione riassegni i blocchi che toccano il \textit{taglio} intelligentemente.
	
% 	\item L'algoritmo è pensato per utilizzare un gran numero di massimi locali come base per descrivere una buona partizione. Anche se nei domini reali, vista la complessità del territorio, i massimi locali sono sempre presenti in un numero abbastanza alto, potrebbe capitare ipoteticamente che questi siano numericamente inferiori al numero di partizioni volute.
% 	Bisognerebbe quindi predisporre un caso limite che permetta di gestire questa configurazione.
% 	In un caso del genere si potrebbe pensare, ad esempio, di suddividere il dominio in base alla  posizione dei blocchi, visto che un basso numero di massimi locali comporta domini equamente distribuiti.
% \end{itemize}

% \subsection*{Le modifiche al modello MPI}
% Come già spiegato precedentemente la versione definitiva del programma \textit{sweGPU} non è stata ancora del tutto sviluppata. Di conseguenza il modello \textit{MPI} dovrà evolversi rispettando le caratteristiche dell'architettura, che si definirà nel corso dello sviluppo della versione distribuita del programma.\\
% Oltre a questa evoluzione il modello \textit{MPI}, dovrà prevedere la possibilità di svolgere comunicazioni considerando anche le possibili multi GPU installate su uno stesso nodo.\\
% Bisognerà inoltre implementare un semplice algoritmo che permetta ai nodi di ricostruire l'ordine delle informazioni ricevute.\\
% Si può pensare ad un semplice ordine nel salvataggio dei dati da inviare. Ad esempio, le informazioni dei bordi dei vari blocchi si possono memorizzare nell'array di dati da inviare, seguendo l'ordine degli identificativi dei blocchi stessi. I dati di uno stesso blocco si possono poi organizzare a seconda della posizione del bordo\footnote{Nord per primo, Est per secondo etc.}.
%\chapter{Obbiettivi}\label{chapter:obbiettivi}
\Blindtext %Dummy Text - remove
%\chapter{Progettazione}\label{chapter:progettazione}
\Blindtext
%\chapter{Risultati}\label{chapter:res}
\Blindtext %Dummy Text - remove
%
%
%%%% Le Conclusioni
\pagestyle{plain}
\chapter*{Conclusione} %Se si cambia il Titolo cambiare anche la riga successiva così che appia corretto nell'conclusione
\addcontentsline{toc}{chapter}{Conclusione} %Per far apparire Introduzione nell'indice (Il nome deve rispecchiare quello del chapter)
Al termine della realizzazione del progetto di stage tutti gli obiettivi descritti nel paragrafo \ref{sec:obiettivi} sono stati raggiunti nei tempi prefissati, 
secondo la ripartizione oraria riportata al paragrafo \ref{sec:ripartizionelavoro}.\\ \\

Il progetto, nella sua fase iniziale, non è risultato di semplice realizzazione in quanto è stato richiesto l'utilizzo di tecnologie e strumenti di cui non si era a conoscenza.\\ \\

Nel corso dello sviluppo ci sono state diverse problematiche:
\begin{itemize}
\item Immagini Docker.
\item Specifiche della Virtual Machine utilizzata.
\end{itemize}

All'inizio si era provato ad utilizzare l'immagine Docker \texttt{ubuntu:latest}, per poi decidere, in modo definitivo, di utilizzare l'immagine ufficiale di Kong Gateway 
per problemi di compatibilità e di raggiungibilità del Gateway.

Ancora, come spiegato nei capitoli precedenti, l'applicazione si trova all'interno di una Virtual Machine sulla quale è stato installato Docker con il relativo container Kong.\\
La problematica era relativa al fatto che in momenti completamente casuali l'applicazione risultava irraggiungibile, e quindi inutilizzabile.\\ 
Andando ad analizzare le statistiche di utilizzo su Microsoft Azure ci si è resi conto che il problema risiedeva nelle risorse (molto limitate) della Virtual Machine, 
che quindi accusava problemi di eccessivo utilizzo di CPU e RAM che implicava l'impossibilità di elaborare le richieste.\\ \\ 

Quindi, parlando di eventuali sviluppi futuri del progetto, sicuramente rientra tra questi il miglioramento delle specifiche della Virtual Machine da utilizzare,
 tenendo conto del fatto che le risorse a disposizione con l'attuale Virtual Machine sono 1 CPU e 512 MB di memoria RAM.\\
Altri sviluppi futuri potrebbero concentrarsi sull'implementazione di un Database più ampio e ottimizzato, andando ad utilizzare ad esempio Microsoft Azure SQL Database, 
il Database integrato nella piattaforma Azure.

%
%%%% La bibliografia
\bibliographystyle{apalike} %{plain} -- Scegliere lo stile preferito
\bibliography{./Bibliografia}
%
%\chapter*{Ringraziamenti} %Se si cambia il Titolo cambiare anche la riga successiva così che appia corretto nell'conclusione
\addcontentsline{toc}{chapter}{Ringraziamenti} %Per far apparire Introduzione nell'indice (Il nome deve rispecchiare quello del chapter)
aaaaaaaaaaaaaaaaaaaaa
%
% Le appendici
\appendix
\chapter{Appendice}
\Blindtext
%
\end{document}